\documentclass[aps,prb,reprint,superscriptaddress,amsfonts,amsmath,amssymb,showpacs,floatfix]{revtex4-1}

\usepackage{graphicx}
\usepackage{multirow}
\usepackage{hyperref}
%\usepackage{txfonts}
\newcommand{\mrm}[1]{\mathrm{#1}}
\newcommand{\mbf}[1]{\mathbf{#1}}
\newcommand{\mbb}[1]{\mathbb{#1}}
\newcommand{\mc}[1]{\mathcal{#1}}
\newcommand{\eref}[1]{(\ref{#1})}
\newcommand{\fref}[1]{Fig.~\ref{#1}}
\newcommand{\tref}[1]{Table~\ref{#1}}
\newcommand{\sref}[1]{Sec.~\ref{#1}}
\newcommand{\nn}{\nonumber}
\newcommand{\rmd}{\mathrm{d}}
\newcommand{\rmi}{\mathrm{i}}
\newcommand{\ra}{\rangle}
\newcommand{\la}{\langle}
\newcommand{\p}[1]{\phantom{#1}}
\newcommand{\rcite}[1]{Ref.~\onlinecite{#1}} % \onlinecite for prb
\newcommand{\prbtext}[1]{#1}
\newcommand{\pratext}[1]{}
\newcommand{\textcitecomma}[1]{\citeauthor{#1},\cite{#1}}
%\newcommand{\textcitecomma}[1]{\textcite{#1},}
\newcommand{\citecomma}[1]{,\cite{#1}}
%\newcommand{\citecomma}[1]{ \cite{#1},}
\newcommand{\citestop}[1]{.\cite{#1}}
%\newcommand{\citestop}[1]{ \cite{#1}.}

\newcommand{\Mop}{M_\alpha^{\p{\alpha}\beta}}
\newcommand{\alphabeta}{_\alpha^{\p{\alpha}\beta}}


\begin{document}

\title{Simulation of anyons with tensor network algorithms}

\author{R. N. C. Pfeifer}
\email[]{pfeifer@physics.uq.edu.au}
\author{P. Corboz}
%\homepage[]{Your web page}
%\thanks{}
%\altaffiliation{}
\affiliation{School of Mathematics and Physics, The University of Queensland, St Lucia, QLD 4072, Australia}
\author{O. Buerschaper}
\author{M. Aguado}
\affiliation{Max-Planck-Institut f\"ur Quantenoptik. Hans-Kopfermann-Stra\ss{}e 1, D-85748 Garching, Germany}
\author{M. Troyer}
\affiliation{Theoretische Physik, ETH Zurich, 8093 Zurich, Switzerland}
\author{G. Vidal}
\affiliation{School of Mathematics and Physics, The University of Queensland, St Lucia, QLD 4072, Australia}

\date{\today}

\begin{abstract}
Interacting systems of anyons pose a unique challenge to condensed matter simulations due to their non-trivial exchange statistics. These systems are of great interest as they have the potential for robust universal quantum computation, but numerical tools for studying them are as yet limited. We show how existing tensor network algorithms may be adapted for use with systems of anyons, and demonstrate this process for the 1-D Multi-scale Entanglement Renormalisation Ansatz (MERA). We apply the MERA to infinite chains of interacting Fibonacci anyons, computing their scaling dimensions and local scaling operators. The scaling dimensions obtained are seen to be in agreement with conformal field theory. The techniques developed are applicable to any tensor network algorithm, and the ability to adapt these ans\"atze for use on anyonic systems opens the door for numerical simulation of large systems of free and interacting anyons in one and two dimensions.
\end{abstract}
%\begin{abstract}
%Interacting systems of anyons pose a unique challenge to condensed matter simulations due to their non-trivial exchange statistics. These systems are of great interest as they have the potential for robust universal quantum computation, but numerical tools for studying them are as yet limited. We demonstrate how the Multi-scale Entanglement Renormalisation Ansatz (MERA) may be adapted for the study of a one-dimensional system of anyons. We apply this ansatz to infinite chains of interacting Fibonacci anyons, computing their scaling dimensions and local scaling operators. The scaling dimensions are seen to be in agreement with conformal field theory. The techniques developed are applicable to any tensor network algorithm, and the ability to adapt these ansatzes for use on anyonic systems opens the door for numerical simulation of large systems of free and interacting anyons in one and two dimensions.
%\end{abstract}
%\\~\\
%\begin{abstract}
%Interacting systems of anyons pose a unique challenge to condensed matter simulations due to their non-trivial exchange statistics. These systems are of great interest as they have the potential for robust universal quantum computation, but numerical tools for studying them are as yet limited. We demonstrate how any existing tensor network algorithm may be adapted for use on a system of anyons, taking as an example the 1-D Multi-scale Entanglement Renormalisation Ansatz (MERA). We apply this ansatz to infinite chains of interacting Fibonacci anyons, computing their scaling dimensions and local scaling operators. The scaling dimensions are seen to be in agreement with conformal field theory. Our techniques open the door to numerical simulation of large systems of free and interacting anyons in one and two dimensions.
%\end{abstract}

% insert suggested PACS numbers in braces on next line
\pacs{05.30.Pr, 73.43.Lp, 02.70.-c, 03.65.Vf}
% insert suggested keywords - APS authors don't need to do this
%\keywords{}

%\maketitle must follow title, authors, abstract, \pacs, and \keywords
\maketitle

% body of paper here - Use proper section commands
% References should be done using the \cite, \ref, and \label commands

\section{Introduction}

The study of anyons offers one of the most exciting challenges in contemporary physics. Anyons are exotic quasiparticles with non-trivial exchange statistics, which makes them difficult to simulate. However, they are of great interest as some species offer the prospect of a highly fault-tolerant form of universal quantum computation\citecomma{kitaev2003,nayak2008} and it has been suggested\prbtext{\cite{xia2004}} that the simplest such species may appear in the fractional quantum Hall state with filling fraction $\nu$ = 12/5\pratext{ \cite{xia2004}}. Despite the current strong interest in the development of practical quantum computing, our ability to study the collective behaviour of systems of anyons remains limited.

The study of interacting systems of anyons using numerical techniques was pioneered by \textcitecomma{feiguin2007} using exact diagonalisation for {1-D} systems of up to 37 anyons, and the Density Matrix Renormalisation Group algorithm (DMRG)\pratext{ }\cite{white1992} for longer chains. Also related is work by \textcitecomma{sierra1997} later extended by \textcitecomma{tatsuaki2000} which applies a variant of DMRG to spin chain models having $SU(2)_k$ symmetry. Some of these models are now known to correspond to $SU(2)_k$ anyon chains\citecomma{trebst2008} and using this mapping these systems may also be studied using the Bethe ansatz\pratext{ }\cite{alcaraz1987} and quantum Monte Carlo\citestop{todo2001}.

However, all of these methods have their limitations. Exact diagonalisation has a computational cost which is exponential in the number of sites, strongly limiting the size of the systems which may be studied. DMRG is capable of studying larger system sizes, but is typically limited to {1-D} or quasi-{1-D} systems (e.g. ladders). Mapping to a spin chain is useful in one dimension but is substantially less practical in two.
There are therefore good reasons to desire
%strong motivation to develop 
a formalism which will allow the application of other tensor network algorithms %to be applied 
to systems of anyons. 
%Tensor networks such as 
Many of these tensor networks, such as 
Projected Entangled Pair States (PEPS)\citecomma{verstraete2004,nishino1998,gu2008,jiang2008,jordan2008}
and the 2-D versions of Tree Tensor Networks (TTN)\pratext{ }\cite{tagliacozzo2009} and of the Multi-scale Entanglement Renormalisation Ansatz (MERA)\pratext{ }\cite{cincio2008,evenbly2009b,evenbly2010}
have been designed specifically to accurately describe two-dimensional systems. 

In one dimension, 
many previously studied systems of interacting anyons display extended critical phases\pratext{ }\prbtext{,}\cite[e.g.][]{feiguin2007,trebst2008}\pratext{,} which are characterised by correlators exhibiting polynomial decay\citestop{difrancesco1997} 
%The study of these systems also stands to benefit from an ansatz which 
Whereas DMRG favours accurate representation of short range correlators at the expense of long-range accuracy, 
%and so there is a clear need for a complementary ansatz which can
% and
%%there is therefore clearly a need for 
%so the study of these systems would stand to benefit from the implementation of
%an %complementary 
%ansatz which can %capable of 
%accurately represent %ing 
%the quasi-long range correlations of these critical anyonic systems.
the {1-D} MERA\pratext{ }\cite{vidal2007,vidal2008a} is ideally suited to this situation as its hierarchical structure naturally encodes the renormalisation group flow at the level of operators and wavefunctions\citecomma{vidal2007,vidal2008a,vidal2010,chen2010} and hence %therefore 
accurately reproduces correlators across a wide range of length scales\citestop{vidal2007,vidal2008a,giovannetti2008,pfeifer2009,evenbly2009} The development of a general formalism for anyonic tensor networks is therefore also advantageous for the study of {1-D} anyonic systems.

%More recently developed tensor networks such as 
%Projected Entangled Pair States (PEPS) \cite{verstraete2004,nishino1998,gu2008,jiang2008,jordan2008},
%and 2-D versions of Tree Tensor Networks (TTN) \cite{tagliacozzo2009} and of the Multi-scale Entanglement Renormalisation Ansatz (MERA) \cite{cincio2008,evenbly2009b,evenbly2010} would facilitate the study of 2-D systems, and even in one 
%
%All of these methods, however, have their limitations. 
%DMRG is typically limited to {1-D} or quasi-{1-D} systems (e.g. ladders), as are approaches which rely on mapping the system of anyons to a spin chain. Exact diagonalisation can be used to study one- or two-dimensional systems, but the computational cost is exponential in the number of sites, strongly constraining the size of the systems which may be studied. For spin systems, many of these problems can be addressed by the application of intrinsically two-dimensional tensor networks such as 
%Projected Entangled Pair States (PEPS) \cite{verstraete2004,nishino1998,gu2008,jiang2008,jordan2008},
%and 2-D versions of Tree Tensor Networks (TTN) \cite{tagliacozzo2009} and of the Multi-scale Entanglement Renormalisation Ansatz (MERA) \cite{cincio2008,evenbly2009b,evenbly2010}.
%There is therefore good reason to desire
%%strong motivation to develop 
%a general formalism which will allow existing tensor networks to be applied to systems of anyons. 
%The development of anyonic tensor networks would also benefit the study of systems of anyons in one dimension, as 
%%The mapping to %the 
%%$SU(2)_k$-symmetric 
%%spin chains is only valid for certain species of anyons, exact diagonalisation has a computational cost which is exponential in the number of sites and strongly limits the size of the systems which may be studied, and
%%DMRG is capable of studying larger system sizes but is typically limited to {1-D} or quasi-{1-D} systems (e.g. ladders), and favours accurate representation of short range correlators at the expense of long-range accuracy. 
%many previously studied systems of interacting anyons display extended critical phases \cite[e.g.][]{feiguin2007,trebst2008}, which are characterised by correlators exhibiting polynomial decay \cite{difrancesco1997}. 
%While DMRG is capable of studying large system sizes, it favours accurate representation of short range correlators at the expense of long-range accuracy, %and
%%there is therefore clearly a need for 
%so the study of these systems would stand to benefit from the implementation of
%an %complementary 
%ansatz which can %capable of 
%accurately represent %ing 
%the quasi-long range correlations of critical anyonic systems.
%The MERA \cite{vidal2007,vidal2008a} is ideally suited to the situation, as its hierarchical structure naturally encodes the renormalisation group flow at the level of operators and wavefunctions \cite{vidal2007,vidal2008a,vidal2010,chen2010}, and it accurately reproduces correlators across a wide range of length scales \cite{vidal2007,vidal2008a,giovannetti2008,pfeifer2009,evenbly2009}.


This paper describes how any tensor network algorithm may be adapted to systems of anyons in one or two dimensions using structures which explicitly implement the quantum group symmetry of the anyon model. As a specific example we demonstrate the construction of the anyonic {1-D} MERA, which we then apply to an infinite chain of interacting Fibonacci anyons at criticality.
%the MERA may be constructed on a {1-D} system of anyons by explicitly implementing the quantum group symmetry of the anyon model. Our approach is, however, 
The approach which we present is %, however, 
completely general, and can be applied to any species of anyons and any tensor network ansatz.
%, including %ansatzes for 2-D systems such as
%Matrix Product States (MPS), PEPS, and {1-D} and 2-D TTN and MERA.
%Projected Entangled Pair States (PEPS) \cite{verstraete2004,nishino1998,gu2008,jiang2008,jordan2008},
%and 2-D versions of Tree Tensor Networks (TTN) \cite{tagliacozzo2009} and of MERA \cite{cincio2008,evenbly2009b,evenbly2010}. 



\section{Anyonic states\label{sec:anyonstates}}

Consider a lattice $\mc{L}_0$ of $n$ sites populated by anyons. In contrast to bosonic and fermionic systems, for many anyon models the total Hilbert space $\mbb{V}_{\mc{L}_0}$ can not be divided into a tensor product of local Hilbert spaces. Instead, a basis is defined by introducing a specific fusion tree (e.g. \fref{fig:anyons}(i)). The fusion tree is always constructed on a linear ordering of %the physical lattice,
anyons, and while the {1-D} lattice naturally exhibits such an ordering, for 2-D lattices some linear ordering %of the lattice sites
must be imposed.
Each line is then labelled with a charge index $a_i$ such that the labels are consistent with the fusion rules of the anyon model,
\begin{equation}
a\times b\rightarrow \sum_{c} N_{ab}^{c} \,c.\label{eq:fusion}
\end{equation}
For anyon types where some entries of the multiplicity tensor $N_{ab}^c$ take values greater than 1, a label $u_i$ is also affixed to the vertex which represents the fusion process to distinguish between the different copies of charge $c$.
The edges of the graph which are connected to a vertex only at their lower end are termed ``leaves'' of the fusion tree, and we will associate these leaves with the charge labels $a_1\ldots a_n$. Different orderings of the leaves on a fusion tree may be interconverted by means of braiding (\fref{fig:anyons}(ii)), and different fusion trees, corresponding to different bases of states, may be interconverted by means of $F$ moves (\fref{fig:anyons}(iii))\citestop{kitaev2006,bonderson2007}
In some situations it may also be useful to associate a further index $b_i$ %(not shown in \fref{fig:anyons}(i)) 
with each of the leaves of the fusion tree. For example, if the leaves are equated with the sites of a physical lattice, then this additional index may be used to enumerate additional non-anyonic degrees of freedom associated with that lattice. %The range of $b_i$ may depend on the value of $a_i$ on the same leaf. 
%Such an index may always be introduced on the leaves of any diagram if required, and for 
For simplicity we will usually leave %it 
these extra indices $b_1\ldots b_n$ implicit, as we have done in \fref{fig:anyons}, 
as they do not directly participate in anyonic manipulations such as $F$ moves and braiding.
% This index behaves exactly like a vertex index $u_i$, 

\begin{figure}
\includegraphics[width=246.0pt]{fusiontree}
\caption{(i)~Example representation of a state $|\psi\ra$ in a fusion tree basis for a system of $6$ anyons. Labels $a_i$ indicate charges associated with edges of the fusion tree graph, and labels $u_i$ are degeneracies associated with vertices. The structure of the tree corresponds to a choice of basis, and does not affect the physical content of the theory.
(ii)~Braiding may be used to change the ordering of the leaves of a fusion tree basis, or to represent anyon exchange.
(iii)~$F$-moves convert between the bases associated with different fusion trees.
%Note: For simplicity the given diagrams assume fusion rules such that $N^c_{ab}$ takes the values 0 or 1. Extension to fusion rules with higher multiplicities is achieved through attaching an additional index to each vertex, which enumerates the different copies of each charge. 
\label{fig:anyons}}
\end{figure}

%In some situations, for example where the leaves are associated with the sites of a physical lattice which has 


Let the total number of charge labels on the fusion tree be given by $m$, where $m\geq n$. 
%For $1\leq i\leq n$ the labels $a_i$ correspond to the charges on the 
%leaves of the fusion tree.
%%lattice sites, represented by the open legs of the tree.
For abelian anyons the fusion rules uniquely constrain all $a_i$ for $i>n$, and provided there are no constraints on the total charge, the total Hilbert space reduces to a product of local Hilbert spaces $\mbb{V}$, such that $\mbb{V}_{\mc{L}_0} = \mbb{V}^{\otimes n}$. For nonabelian anyons, additional degrees of freedom arise because some fusion rules admit multiple outcomes, permitting certain $a_i\ (i>n)$ to take on multiple values while remaining consistent with the fusion rules, and the resulting Hilbert space does not necessarily admit a tensor product structure.

We will now associate a parameter $\nu_{i,a_i}$ with each charge on the fusion tree, which we will term the degeneracy. This parameter corresponds to the number of possible fusion processes by which charge $a_i$ may be obtained at location $i$. Where charge $a_k$ arises from the fusion of charges $a_i$ and $a_j$, then $\nu_{k,a_k}$ will satisfy
\begin{equation}
\nu_{k,a_k} = \sum_{a_i,a_j} \nu_{i,a_i}\nu_{j,a_j} N^{a_k}_{a_i a_j}.\label{eq:compounddegens}
\end{equation}
For systems where the only degrees of freedom are anyonic, degeneracies on the physical lattice $\mc{L}_0$ (i.e. $\nu_{i,a_i}$, $1\leq i\leq n$) will take values of 0 or 1 depending on whether a charge $a_i$ is permitted on lattice site $i$. Higher values of $\nu_{i,a_i}$ may be used on the physical lattice 
if there is also a need to represent additional 
non-anyonic degrees of freedom, %as would be 
enumerated by indices $b_1\ldots b_n$.

Up to this point we have parameterised our Hilbert space in terms of explicit labellings of the fusion tree. We now adopt a different approach: Consider an edge $i$ of the fusion tree which is not a ``leaf''. %does not connect directly to the physical lattice. 
As well as labelling this edge with a charge $a_i$ we may introduce a second index $\mu_i$, running from 1 to $\nu_{i,a_i}$. Each pair of values $\{a_i,\mu_i\}$ may be associated with a unique charge labelling for the portion of the fusion tree 
from edge $i$ out to the leaves,
%which connects edge $i$ to the physical lattice,
with these labellings being compatible with the fusion rules in the presence of a charge of $a_i$ on site $i$ (for an illustration of this, see \fref{fig:exampleedges}). 
\begin{figure}
\includegraphics[width=246.0pt]{exampleedges}
\caption{%On this fusion tree, the physical lattice sites are associated with the open edges of the fusion tree labelled 
The leaves of this fusion tree carry the charge labels
$a_1$ to $a_6$. %All other edges $a_i$, $i>6$ are therefore edges which do not connect directly to the physical lattice. 
An edge which is not a leaf, labelled with charge $a_{10}$, is indicated by the large grey arrow. 
%For this edge, the portion of the fusion tree which connects it to the physical lattice 
The portion of the fusion tree extending from edge $a_{10}$ out to the leaves
is indicated by the grey ellipse. If a degeneracy index $\mu_{10}$ is associated with charge $a_{10}$, then for a given value of $a_{10}$, index $\mu_{10}$ will 
%the edge labelled $a_10$ (and indicated by the large dashed arrow) is an example of an edge which does not connect to the physical lattice. The portion of the fusion tree which connects edge $a_10$ to the physical lattice is surrounded by the grey ellipse. W
%Indicated on this fusion tree are (i) an example of an edge which does not connect directly to the physical lattice (i.e. an edge of the graph labelled by a charge $a_i$, $i>n$), and (ii) the portion of the fusion tree which connects this edge to the physical lattice. For a given value of the charge $a_i$, index $\mu_i$ 
enumerate all compatible labellings of the highlighted portion of the fusion tree.\label{fig:exampleedges}}
\end{figure}%
Provided we know the structure of the fusion tree above $i$ and have a systematic means of associating labellings of that portion of the tree with values of $\mu_i$, then in lieu of stating the values of all $a_j$ for edges $j$ involved in that portion of the tree, we may simply specify the value of the degeneracy index $\mu_i$. 
In this way we may specify an entire state in the form
\begin{equation}
|\psi\ra = \sum_{\mu_{m}} c_{a_{m}\mu_{m}} |a_{m},\mu_{m}\ra\label{eq:statepsi_pre}
\end{equation}
where $a_{m}$ is the total charge obtained on fusing all the anyons. %on the physical lattice. 
The index $\mu_{m}$, which is the degeneracy index associated with the total charge of the fusion tree, may be understood as systematically enumerating all possible labellings of the entire fusion tree including charge labels, vertex labels, and any labels associated with additional non-anyonic degrees of freedom. % $b_1\ldots b_n$. %on the physical lattice. 
For an example, see \fref{fig:exampleenum}.
\begin{figure}
\includegraphics[width=246.0pt]{exampleenum}
\caption{Example enumeration of states according to $a_m$ and $\mu_m$ for a fusion tree describing %the fusion of 
four Fibonacci anyons. The Fibonacci anyon model has one non-vacuum charge label ($\tau$) and one non-trivial fusion rule, $\tau\times\tau\rightarrow 1+\tau$. Because the charges $1$ and $\tau$ are both self-dual, no arrows are required on diagrammatic representations of Fibonacci anyon fusion trees.\label{fig:exampleenum}}
\end{figure}%
Note that for a given edge $i$, the value of the degeneracy $\nu_{i,a_i}$ may vary with the charge $a_i$ and consequently the range of the degeneracy index $\mu_{m}$ in Eq.~\eref{eq:statepsi_pre} is dependent on the value of the charge $a_{m}$.

The notation of Eq.~\eref{eq:statepsi_pre} should be contrasted with that of
%Contrast this with the notation of 
\fref{fig:anyons}(i). In the latter, the number of indices on $c$ depends upon the number of charge labels on the fusion tree, whereas in the former, the tensor describing the state is always indexed by just one pair of labels---charge and degeneracy---which will prove advantageous in constructing a tensor network formalism for systems of anyons.

We now choose to restrict our attention to systems having the identity charge. We may do this without loss of generality as a state on $n$ lattice sites with a total charge $a_{m}$ may always be equivalently represented by a state on $n+1$ lattice sites whose total charge is the identity, with a charge $\overline{a_{m}}$ on lattice site $n+1$. This additional charge annihilates the total charge $a_{m}$ of sites $1\ldots n$ to give the vacuum. The expression for $|\psi\ra$ then becomes
\begin{equation}
|\psi\ra = \sum_{\mu_{m'}} c_{1\mu_{m'}} |1,\mu_{m'}\ra\label{eq:statepsi}
\end{equation}
where $\mu_{m'}$ ranges from 1 to the dimension of the Hilbert space of the system of $n$ sites with total charge $a_m$.
Consequently we may represent the state $|\psi\ra$ of a system of anyons by means of the vector $c_{1\mu_{m'}}$. For simplicity of notation, we will take greek indices from the beginning of the alphabet to correspond to pairs of indices $\{a_i,\mu_i\}$ consisting of a charge index and the associated degeneracy index. The vector $c_{1\mu_{m'}}$ will therefore be denoted simply $c^\alpha$, with the understanding that in this case 
%$c^\alpha$ is nonzero only when 
the charge component $a_{m'}$ of multi-index $\alpha$ takes 
only
the value 1. (Multi-index $\alpha$ is raised as we will shortly introduce a diagrammatic formalism in which vector $c$ is represented by an object with a single upward-going leg. In this formalism, upward- and downward-going legs may be associated with upper and lower multi-indices respectively.)

%It is worth noting that for anyon models where multiplicities $N^c_{ab}$ may be greater than 1, such as $SU(3)_k$, the degeneracy indices $\mu_i$ absorb the vertex indices which distinguish between multiple copies of a particular charge.


\section{Anyonic operators}

We will divide our consideration of anyonic operators into two parts. First we shall consider operators which map a state on some Hilbert space $\mc{H}$ into another state on the same Hilbert space. When applied to a state represented by $c^\alpha$, such an operator leaves the degeneracies of the charges in multi-index $\alpha$ unchanged. We will therefore call these
\emph{degeneracy-preserving} anyonic operators.
Then we will consider those operators which map a state on some Hilbert space $\mc{H}$ into a state on some other Hilbert space $\mc{H}'$. These operators may represent processes which modify the environment, for example by adding or removing lattice sites, and also play an important part in anyonic tensor networks, for instance taking the role of isometries in the TTN and MERA. 
%These operators may also be written in the form $\Mop$ but do not satisfy the additional constraints on degeneracies described above. 
%We will therefore term these 
As these operators can change the degeneracies of charges in a multi-index $\alpha$, we will call them \emph{degeneracy-changing} anyonic operators.
%However, these operators cannot be represented in the manner of \fref{fig:operators}(i), so we shall defer further discussion of them until \sref{sec:degenexp}, after developing a more powerful diagrammatic notation.
More generally, the degeneracy-preserving anyonic operators may be considered a subclass of the degeneracy-changing anyonic operators for which $\mc{H}=\mc{H}'$.

\subsection{Degeneracy preserving anyonic operators\label{sec:anyonops}}

%The first class of operators which we are interested in consists of the operators mapping 
We begin with those operators which map
states on some Hilbert space $\mc{H}$ into other states on the same Hilbert space $\mc{H}$. 
Examples of these operators include Hamiltonians, reduced density matrices, and unitary transformations such as the disentanglers of the MERA.

First, we introduce splitting trees. The space of splitting trees is dual to the space of fusion trees. While the space of fusion trees consists of labelled directed graphs whose number of branches increases monotonically when read from bottom to top, the space of splitting trees consists of labelled directed graphs whose number of branches increases monotonically when read from top to bottom. 
An inner product is defined by connecting the leaves of fusion and splitting trees which have equivalent linear orderings of the leaves (%using 
braiding first %to ensure this 
if necessary), 
then eliminating all loops as per \fref{fig:operators}(i), with $F$ moves performed as required. % to eliminate all loops.
%For fusion and splitting trees with equivalent orderings of leaves, the inner product may be 
%
%according to \fref{fig:operators}(i).
\begin{figure}
\includegraphics[width=246.0pt]{operators}
\caption{
(i) Loops are eliminated by replacing them 
%To evaluate inner product between a fusion and a splitting tree, the leaves of the trees are connected, with braiding if required, then $F$ moves are performed and loops are eliminated. When a loop is eliminated, it is replaced 
with an equivalent numerical factor determined by the normalisation convention. 
%The removal of a loop is illustrated here, with the numerical factor corresponding 
The factor given here corresponds
to the diagrammatic isotopy convention employed in \protect{\rcite{bonderson2007}}.
(ii) Definition of a simple two-site anyonic operator. 
(iii) Application of an operator to a state is performed by connecting the diagrams' free legs. By performing $F$ moves and eliminating loops (and in more complex examples, also braiding) it is possible to obtain an expression for the resulting state in the original basis. 
%As in \protect{\fref{fig:anyons}}, the diagrams presented assume fusion rules with multiplicities of 0 or 1, but may be extended to larger multiplicities by the incorporation of vertex indices.
\label{fig:operators}}
\end{figure}

Anyonic operators
may always be written as a sum over fusion and splitting trees, such as the two-site operator $\hat M$
shown in \fref{fig:operators}(ii), and for degeneracy-preserving anyonic operators it is %also 
always possible to choose the splitting tree to be the adjoint of the fusion tree.
%and it is always possible to rearrange these trees such that one is the adjoint of the other.
To apply an operator %of this type 
to a state the two corresponding diagrammatic representations are connected as shown in \fref{fig:operators}(iii), and closed loops may be eliminated as shown in \fref{fig:operators}(i). 
%The coefficient associated with elimination of a loop arises from the isotopy invariance convention described in \rcite{bonderson2007}, and 
Sequences of $F$ moves, braiding, and loop eliminations may be performed until the diagram has been reduced once more to a fusion tree without loops on a lattice of $n$ sites.

Much as the state of an anyonic system may be represented by a vector $c^\alpha$, anyonic operators may be represented by a matrix $\Mop$. Each value of $\alpha$ corresponds to a pair $\{a_i,\mu_i\}$ where $a_i$ is a possible charge of the central edge of the operator diagram (e.g. $a_3$ in \fref{fig:operators}(ii)), and $\mu_i$ is a value of the degeneracy index associated with charge $a_i$. We will denote the degeneracy of $a_i$ by $\nu_{a_i}$. Similarly, values of $\beta$ correspond to pairs $\{a_j,\mu_j\}$ where $a_j$ has degeneracy $\nu_{a_j}$.
For degeneracy-preserving anyonic operators the charge indices
$a_i$ and $a_j$ necessarily take on the same range of values, and $\nu_{a_i}=\nu_{a_j}$ when $a_i=a_j$. The values of $\nu_{a_i}$ may equivalently be calculated from either the fusion tree making up the top half or the splitting tree making up the bottom half of the operator diagram.

A well-defined anyonic operator $\hat M$ must respect the (quantum) symmetry group of the anyon model, and consequently all entries in $\Mop$ for which $a_i\not=a_j$ will be zero. However, in contrast with $c^\alpha$ we do not require that $a_i=a_j=1$. When $\hat M$ is a degeneracy-preserving operator, matrix $\Mop$ is therefore a square matrix of side length
\begin{equation}
\ell_M=\sum_{a_i} \nu_{a_i},
\end{equation}
which may be organised to exhibit a structure which is block diagonal in the charge indices $a_i$ and $a_j$, and for which the blocks are also square. 
%these blocks will be square.
As an example consider \fref{fig:exampleoperator}, which shows an operator acting on four Fibonacci anyons. An example matrix $\Mop$ for an operator of this form is given in \tref{tab:examplematrix}, from which the entries of $M_{abcde}$ can be reconstructed, e.g. $M_{\tau 1 \tau 1 \tau}=3$.
\begin{figure}
\includegraphics[width=246.0pt]{exampleoperator}
\caption{An operator acting on four Fibonacci anyons. The values of the coefficients $M_{abcde}$ may be specified as a block-diagonal matrix $M_\alpha^{~\beta}$, for example as in \protect{\tref{tab:examplematrix}}.\label{fig:exampleoperator}}
\end{figure}%

\begin{table}
\begin{equation*}
\Mop=\quad\begin{array}{cc|ccccccc}
&&&\multicolumn{5}{c}{a_j,\mu_j}&\\
&&&1,1&1,2&\tau,1&\tau,2&\tau,3&\\
\hline
\multirow{5}{*}{$a_i$,$\mu_i$}&1,1&
\multirow{5}{*}{$\left(\begin{array}{c}\!\\ \!\\ \!\\ \!\\ \!\end{array}\right.$}& 1&0.5&0&0&0 &
\multirow{5}{*}{$\left.\begin{array}{c}\!\\ \!\\ \!\\ \!\\ \!\end{array}\right)$}\\
&1,2&& 0.5&1&0&0&0 &\\
&\tau,1&& 0&0&1&2&-1 &\\
&\tau,2&& 0&0&2&3&-1 &\\
&\tau,3&& 0&0&1&1&1 &
\end{array}
\end{equation*}
~

\begin{tabular}{|cc|ccccc|}
\hline
$a_i$&$\mu_i$&~&$a$&$b$&$c$&\\
\hline\hline
%\multirow{2}{*}{1}&1&1&$\tau$&1\\
1&1&&1&$\tau$&1&\\
1&2&&$\tau$&$\tau$&1&\\
$\tau$&1&&1&$\tau$&$\tau$&\\
$\tau$&2&&$\tau$&1&$\tau$&\\
$\tau$&3&&$\tau$&$\tau$&$\tau$&\\
\hline
\end{tabular}
~~~~
\begin{tabular}{|cc|ccccc|}
\hline
$a_j$&$\mu_j$&~&$e$&$d$&$c$&\\
\hline\hline
1&1&&1&$\tau$&1&\\
1&2&&$\tau$&$\tau$&1&\\
$\tau$&1&&1&$\tau$&$\tau$&\\
$\tau$&2&&$\tau$&1&$\tau$&\\
$\tau$&3&&$\tau$&$\tau$&$\tau$&\\
\hline
\end{tabular}
\caption{Matrix representation $M_\alpha^{~\beta}$ for an example operator of the form shown in \protect{\fref{fig:exampleoperator}}. Multi-index $\alpha$ corresponds to index pair $\{a_i,\mu_i\}$ and multi-index $\beta$ corresponds to pair $\{a_j,\mu_j\}$. Subject to an appropriate ordering convention for $\mu_i$ and $\mu_j$, these indices may be related to the fusion tree labels $a,b,c,d,e$ of \protect{\fref{fig:exampleoperator}} as shown. Note that as $c$ is the charge on the central leg of \protect{\fref{fig:exampleoperator}}, all nonzero entries of $M_\alpha^{~\beta}$ satisfy $a_i=a_j=c$.\label{tab:examplematrix}}
\end{table}

% A defining property of this class of anyonic operators is that they map a state on some Hilbert space $\mathcal{H}$ into another state in the same Hilbert space. It is also possible to define anyonic operators which map states from one Hilbert space $\mc{H}$ to another Hilbert space $\mc{H}'$. These operators may represent processes which modify the environment, for example by adding or removing lattice sites. They also play an important part in anyonic tensor networks, for instance taking the role of isometries in the TTN and MERA. However, these operators cannot be represented in the manner of \fref{fig:operators}(ii), so we shall defer further discussion of them until \sref{sec:degenexp}, after developing a more powerful diagrammatic notation. 


\subsection{Degeneracy changing anyonic operators\label{sec:degenexp}}

We now introduce the second class of anyonic operators, which map states in some Hilbert space $\mc{H}$ into some other Hilbert space $\mc{H}'$. These operators may reduce or increase the degeneracy of any charge present in the spaces on which they act, and may even project out entire charge sectors by setting their degeneracy to zero.
% which modify the size of the degeneracy space. 
When these operators are written in the conventional notation of \fref{fig:operators}, the fusion and splitting trees will not be identical. Further, we may choose to allow combinations of degeneracies which do not naturally admit complete decomposition into individual anyons. For example, a degeneracy-changing operator may map a state on five Fibonacci anyons (having total degeneracies $\nu_1=3$, $\nu_\tau=5$) into a state having degeneracies $\nu_1=2$, $\nu_\tau=2$. As these degeneracies do not admit decomposition into an integer number of nondegenerate anyons, it is necessary %instead 
to associate an index $u_i$ with the single open leg of the fusion tree. This index behaves identically to the vertex indices $u_i$ of \fref{fig:anyons}, serving to enumerate the different copies of each individual charge, and as with the vertex indices of \fref{fig:anyons}, it is absorbed into the degeneracy index $\mu_i$.

As a further example, a state having degeneracies $\nu_1=4$, $\nu_\tau=4$ could be associated with a fusion tree having either one leg, or two legs each with degeneracies $\nu_1=0$, $\nu_\tau=2$. Again, indices $u_i$ would have to be associated with each open leg.

%Fibonacci anyons do not require indices $u_i$ to be associated with the vertices of their diagrams, but since this state does not admit decomposition into any integer number of identical nondegenerate anyons, it will require the association of an index $u_i$ with each of the open legs of the $\nu_1=4$, $\nu_\tau=4$ fusion or splitting tree

%. For this example, possible choices of tree include a single leg with degeneracies $\nu_1=4$, $\nu_\tau=4$, or two legs each with degeneracy $\nu_1=0$, $\nu_\tau=2$. 

%Similarly, some states may only admit partial decomposition, e.g. a state having $\nu_1=4$, $\nu_\tau=4$ may be considered the fusion of two sites each having $\nu_1=0$, $\nu_\tau=2$, and again a degeneracy index must be associated with each of the open legs to distinguish between the two different copies of $\tau$.

%As with degeneracy-preserving anyonic operators, 
%If the fusion or splitting space cannot be written as the fusion of nondegenerate charges, it may also be necessary to associate degeneracy indices with the open legs of either or both of the fusion and splitting trees.
Matrix representations of degeneracy-changing anyonic operators may also %again 
be constructed, and when they %matrix representations of degeneracy-changing anyonic operators 
are written in block diagonal form, the matrices and their blocks may be rectangular rather than square. Degeneracy-changing anyonic operators therefore 
represent a generalisation of the degeneracy-preserving anyonic operators discussed in \sref{sec:anyonops}. It is worth noting that the presence of indices $u_i$ on the open legs of the fusion or splitting trees of an operator do not automatically imply that it is a degeneracy-changing anyonic operator: The defining characteristic of a degeneracy-preserving anyonic operator is that it maps a state in a Hilbert space $\mc{H}$ into a state in the same Hilbert space $\mc{H}$, and consequently both the matrix as a whole and all of its blocks are square. Thus a degeneracy-preserving anyonic operator may act on states having additional indices $u_i$ on their open legs, and the resulting state may be expressed in the form of the same fusion tree, with the same additional indices on the open legs.

% As described in  \sref{sec:anyonstates}, any additional indices which had to be associated with the open legs of the fusion tree 
%Such operators may reduce or increase the degeneracy of any charge already present in the space on which they act, and may even project out entire charge sectors if these sectors have zero degeneracy on one index. 

%These tensors may again be associated with diagrams of the form shown in \fref{fig:blobtensors}(i), consisting of the matrix representation of the operator (represented by the central grey object), and unlabelled fusion trees. In particular, in the manner explained in \sref{sec:anyonstates}, the degeneracy components of the multi-indices automatically incorporate any additional indices which had to be associated with the open legs of the trees.

%Such tensors are again associated with diagrams of the form shown in \fref{fig:blobtensors}(i), but may not always be written in the form of \fref{fig:operators}(ii) as the charges and degeneracies on either side of the central matrix object in \fref{fig:blobtensors} do not in general coincide.
%In addition, the fusion and splitting trees above and below the central object may differ in their degeneracies, admissible charges, and even in their structure.
Operators which change degeneracies may %of this sort may %be used to 
represent physical processes which change the accessible Hilbert space of a system. %, or in the construction of an efficient tensor network representation of a state. 
%Their fusion and splitting trees may be manipulated, and they may be raised and combined with other operators, in exactly the same manner as the simpler operators described in \sref{sec:anyonops}, of which they are a useful generalisation. 
%These objects 
As we will see in \sref{sec:MERAconstr}, 
they may also be used in tensor network algorithms as part of an efficient representation of particular states or subspaces of a Hilbert space, for example the ground state or the low energy sector of a local Hamiltonian.
%the otherwise exponentially large Hilbert space, for example the ground state or low energy sector of a local Hamiltonian.
%Examples of this type of operator will be seen in \sref{sec:MERAconstr}. 

This distinction between degeneracy-changing and degeneracy-preserving anyonic operators is %perhaps most 
clearly seen with a simple example. Let $|\psi\ra$ be a state on six Fibonacci anyons. This state can be parameterised by a vector $c^\alpha$, which has five components. We now define two projection operators, $\hat P^{(1)}$ and $\hat P^{(2)}$ (\fref{fig:projectionops}), each of which acts on the fusion space of anyons $\tau_1$ and $\tau_2$. 
\begin{figure}
%\begin{center}
\includegraphics[width=246.0pt]{projectionops}
\caption{Diagrammatic representation of operators $\hat P^{(1)}$ \eref{eq:P1} and $\hat P^{(2)}$ \eref{eq:P2}. The charges on all leaves are non-degenerate.\label{fig:projectionops}}
%\end{center}
\end{figure}%
Operator $\hat P^{(1)}$ is degeneracy-preserving, and projects $c^\alpha$ into the subspace in which anyons $\tau_1$ and $\tau_2$ fuse to the identity. Its matrix representation is
\begin{equation}
P_{\p{(1)}\,\alpha}^{(1)\,\p{\alpha}\beta} =
%\begin{array}{ccccc}
%&&1&\tau&\\
%1&\multirow{2}{*}{$\left(\begin{array}{c}\!\\\!\end{array}\right.$}&
%1&0&
%\multirow{2}{*}{$\left.\begin{array}{c}\!\\\!\end{array}\right)$}\\
%\tau&&0&0&
%\end{array}
\left(
\begin{array}{cc}
1&0\\0&0
\end{array}
\right)\label{eq:P1}
\end{equation}
where the first value of each multi-index corresponds to a charge of 1, and the second to a charge of $\tau$.
Operator $\hat P^{(2)}$ performs the same projection, but is degeneracy-changing. Its matrix representation is written
\begin{equation}
P_{\p{(2)}\,\alpha}^{(2)\p{\,\alpha}\beta}=(~1~0~).%\left(\begin{array}{cc}1&0\end{array}\right).
\label{eq:P2}
\end{equation}
%The tree structures associated with these operators are shown in \fref{fig:projectionops}.
Both operators perform equivalent projections, in the sense that %and hence 
\begin{equation}
\la\psi|\hat P^{(1)\dagger}\hat P^{(1)}|\psi\ra = \la\psi|\hat P^{(2)\dagger}\hat P^{(2)}|\psi\ra.
\end{equation} 
When $\hat P^{(1)}$ acts on $|\psi\ra$ it leaves the Hilbert space unchanged, and hence the vector $c'^\alpha$ describing state $|\psi'\ra=\hat P^{(1)}|\psi\ra$ is once again a five-component vector, although in an appropriate basis some components will now necessarily be zero. In contrast $\hat P^{(2)}$ explicitly reduces the dimension of the Hilbert space, and the vector $c''^\alpha$ describing state $|\psi''\ra=\hat P^{(2)}|\psi\ra$ is of length two, describing a fusion tree on only four Fibonacci anyons (as both $\tau_1$ and $\tau_2$ have been eliminated). One consequence of this distinction is that while $(\hat P^{(1)})^2=\hat P^{(1)}$, the value of $(\hat P^{(2)})^2$ is undefined.

\section{Anyonic tensor networks}

\subsection{Diagrammatic notation\label{sec:tensordiagrammatic}}

The diagrammatic notation conventionally employed in the study of anyonic systems, and used here in Figs.~\ref{fig:anyons} and \ref{fig:operators}, is well suited to the complete description of anyonic systems, as it provides a physically meaningful depiction of the entire Hilbert space. 
However, the number of parameters required for such a description grows exponentially in the system size, and because it is necessary to explicitly assign every index to a specific charge or degeneracy, specification of a tensor network rapidly becomes inconveniently verbose (for example see \fref{fig:operators}(iii)). %, can become inconveniently verbose.
%The motivation behind tensor networks is to approximately represent certain portions of this Hilbert space, corresponding for instance to the low energy sector of a local Hamiltonian, in a more efficient manner and we must first develop a notation which allows us to do that.
%We will now introduce a notation which permits us to represent tensor networks in a more efficient manner.

In the preceding Sections, we developed techniques whereby anyonic states and operators could be %efficiently 
represented as vectors and matrices, bearing only one or two multi-indices apiece. We now introduce the graphical notation which complements this description, and in which we will formulate anyonic tensor networks. \fref{fig:blobtensors}(i) gives the graphical representations of a state $|\psi\ra$ associated with a vector $c^\alpha$, and of an operator $\hat M$ associated with a matrix $\Mop$.
\begin{figure}
\includegraphics[width=246.0pt]{blobtensors}
\caption{(i) Diagrammatic representation of a state $|\psi\ra$ and two-site operator $\hat M$ expressed in terms of degeneracy indices. (ii) Application of $\hat M$ to state $|\psi\ra$. Grey shapes represent tensors with charge and degeneracy multi-indices, with each leg of the shape corresponding to one charge and degeneracy index pair. These diagrams represent the same state, operator, and process as \protect{\fref{fig:anyons}}(i) and \protect{\fref{fig:operators}}(ii)-(iii).\label{fig:blobtensors}}
\end{figure}
The circle marked $c$ corresponds to the vector $c^\alpha$, and the circle marked $M$ corresponds to the matrix $\Mop$. In general, grey circles correspond to tensors, and the number of legs on the circle corresponds to the number of multi-indices on the associated tensor. Each multi-index is also associated with a fusion or splitting tree structure, which is specified graphically. For reasons to be discussed shortly, we will require that no tensor ever have more than three multi-indices.
As the legs of the grey shapes are each associated with a multi-index, they carry both degeneracy and charge indices. Consequently it is not necessary to explicitly assign labels to the fusion/splitting trees, as these labellings are contained implicitly in the degeneracy index (for example see \tref{tab:examplematrix}, where specifying the values of $\{a_i,\mu_i\}$ and $\{a_j,\mu_j\}$ is equivalent to fully labelling the fusion and splitting trees of \fref{fig:exampleoperator}). 

The fusion or splitting tree associated with a particular multi-index may be manipulated in the usual way by means of braids and $F$ moves, recalling that each component of the tensor is associated with a particular labelling of the fusion and splitting trees via the corresponding values of the multi-indices. %Braids and $F$ moves
Manipulations performed upon a particular tree thus generate unitary matrices which act upon the multi-index that corresponds to the labellings of that particular tree.

The application of an operator to a state is, unsurprisingly, performed by connecting the appropriate diagrams, as shown in \fref{fig:blobtensors}(ii). For operators of the type discussed in \sref{sec:anyonops}, the outcome is necessarily a new state in the same Hilbert space, which consequently can be described by a new state vector $c'^\alpha$, as shown. However, in general an operator $\hat M$ will not act on the entire Hilbert space of the system, and so will be described by a tensor constructed on the fusion space of some subset of lattice sites, and not on the system as a whole. Operator $\hat M$ acting on state $|\psi\ra$ in \fref{fig:blobtensors}(ii) is an example of this. Because $c^\alpha$ describes a six-site system but $\Mop$ is constructed on the fusion space of two sites, 
the multi-indices of $\Mop$ span a significantly smaller Hilbert space than that of $c^\alpha$ and we cannot simply write 
\begin{equation}
c'^\beta = c^{\alpha}M_\alpha^{\phantom{\alpha}\beta}
\end{equation}
(using Einstein notation, where repeated multi-indices are assumed to be summed).
Instead, we must understand how to expand the matrix representation of an operator on some number of sites $x$, to obtain its matrix representation as an operator on $x'$ sites, where $x'>x$. 

%\section{Anyonic tensor networks}

%When constructing an anyonic tensor network, it will often be necessary to act 

\subsection{Site expansion of anyonic operators\label{sec:siteexpanyop}}

The multiplicity tensor $N^{c}_{ab}$ describes the fusion of two charges without degeneracies. It is easily extended to incorporate degeneracies of the charges, 
and we will denote this expanded multiplicity tensor $\tilde N^\gamma_{\alpha\beta u}$ 
where %values of 
multi-indices $\alpha$, $\beta$, and $\gamma$ are associated with the pairs 
$\{a,\mu_a\}$, $\{b,\mu_b\}$, and $\{c,\mu_c\}$ respectively, and for given values of $\alpha$, $\beta$, and $\gamma$, $u$ runs from 1 to $N^c_{ab}$. 
The degeneracies associated with charges $a$, $b$, and $c$ are denoted $\nu_{a}$, $\nu_{b}$, and $\nu_c$ respectively. As with $\mu_a$, $\mu_b$, and $\mu_c$, there is an implicit additional index on each degeneracy $\nu_x$ representing the edge of the tree on which %the associated 
charge $x$ resides. 
The values of $\nu_a$ and $\nu_b$ may be chosen arbitrarily (for example, $\nu_{a}|_{a=1}$ may differ from $\nu_{b}|_{b=1}$), but the degeneracies associated with the values of $c$ must satisfy
\begin{equation}
\nu_{c}=\sum_{a,b} \nu_{a}\nu_{b}N^{c}_{ab}
\end{equation}
in accordance with Eq.~\eref{eq:compounddegens}. When this constraint is satisfied, 
every quadruplet of indices $\{a,\mu_a,b,\mu_b\}$ corresponding to a unique pair of choices for $\alpha$ and $\beta$ may be associated with $N^c_{ab}$ distinct
pairs of indices $\{c,\mu_c\}$ for each $c\in a\times b$. These pairs $\{c,\mu_c\}$ are enumerated by the additional index $u$. % running from 1 to $N^c_{ab}$, 
This defines a 1:1 mapping between sets of values on $\{a,\mu_a,b,\mu_b,u\}$ and pairs $\{c,\mu_c\}$,
%
%corresponding to a unique choice of $\gamma$ may be associated with $N^c_{ab}$ 
%distinct quadruplets
%%set 
%of
%%five
%indices $\{a,\mu_a,b,\mu_b\}$ %, each corresponding to a unique %pair of choices for $\alpha$ and $\beta$, and 
%satisfying $c\in a\times b$. %For groups where all entries in $N^c_{ab}$ are zero or one, this mapping is one-to-one.
%Having made this mapping between values of $\gamma$ and %pairs
%sets of values on $u$, $\alpha$, and $\beta$, 
and we set the corresponding entries in $\tilde N^\gamma_{\alpha\beta u}$ to 1, with all other entries being zero.
%take the value $(N^c_{ab})^{-\frac{1}{2}}$, with all other entries in $\tilde N^\gamma_{\alpha\beta u}$ being zero. %For groups with all multiplicities $N^c_{ab}\leq 1$, $\tilde N^c_{ab}$ satisfies $N^\gamma_{\alpha\beta}\tilde N^{\dagger\alpha\beta}_{\p\dagger\epsilon}=\delta_\epsilon^{\p\epsilon\gamma}$. 
A simple example is given in \tref{tab:exampleN}.
\begin{table}
\begin{tabular}{c|c}
Pair & Assigned pentuplet \\
\{$c,\mu_c$\} & \{$a,\mu_a,b,\mu_b,u$\} \\
\hline
$1,~1$ & $1,~1,~1,~1,~1$ \\
$1,~2$ & $\tau,~1,~\tau,~1,~1$ \\
$1,~3$ & $\tau,~1,~\tau,~2,~1$ \\
$\tau,~1$ & $1,~1,~\tau,~1,~1$ \\
$\tau,~2$ & $1,~1,~\tau,~2,~1$ \\
$\tau,~3$ & $\tau,~1,~1,~1,~1$ \\
$\tau,~4$ & $\tau,~1,~\tau,~1,~1$ \\
$\tau,~5$ & $\tau,~1,~\tau,~2,~1$
\end{tabular}
\caption{Construction of $\tilde N^\gamma_{\alpha\beta u}$ for a fusion vertex for Fibonacci anyons.
In this example $a$ may take charges $1$ and $\tau$ each with degeneracy 1, and $b$ may take charges $1$ and $\tau$ with degeneracies 1 and 2 respectively. By Eq.~\protect{\eref{eq:compounddegens}}, charge $c$ may therefore take values $1$ and $\tau$ with degeneracies 3 and 5 respectively.
A correspondence between the values of multi-index $\gamma$ and of multi-indices $\alpha$ and $\beta$ is established in some systematic manner, with each assignation satisfying $c\in a\times b$, and for Fibonacci anyons the index $u$ is trivial as all multiplicities $N^c_{ab}$ are zero or one.
%For anyon models with multiplicities greater than 1, index $u$ ensures that the correspondence between $\{c,\mu_c\}$ and $\{a,\mu_a,b,\mu_b,u\}$ is one-to-one.
An example %of such an 
assignation is shown in the table. The corresponding entries of $\tilde N$ are then set to %$\sqrt{N^c_{ab}}=1$
1, with all other entries zero. For example, the fourth row indicates that $\tilde N^{(\tau,1)}_{(1,1)(\tau,1)1}=1$.\label{tab:exampleN}}
\end{table}

%The following example may clarify the relationship between these two formulations: Consider the fusion rule $q\times q\rightarrow \mbb{I}+2q+\ldots$. An operator acts on two sites, each carrying charge $q$, and so maps $q\times q$ into $q\times q$. This operator may be written as a block-diagonal matrix. In the sector acting on states with total charge $q$, we may write this block either as $B_{\mu_a}^{\p{\mu_a}\mu_b}\otimes[(\mc{R}_q\oplus\mc{R}_q)\otimes (\mc{R}_q\oplus\mc{R}_q)]$, or as $B_{\mu_a}^{\p{\mu_a}\mu_b}\otimes\mc{R}_{

%If $\tilde N$ is decomposed into a spin network and a degeneracy space and for some $\{a,b,c\}$ we have $N^c_{ab}>1$, then instead of associating $N^c_{ab}$ identical copies of a tensor $\mbb{D}^c_{ab}$ on the degeneracy space with the $N^c_{ab}$ different copies of charge $c$, we associate a single copy of $\mbb{D}^c_{ab}$ with $c\oplus

By virtue of their derivation from $N^{c}_{ab}$, the object $\tilde N^\gamma_{\alpha\beta u}$ and its conjugate $\tilde N^{\dagger\alpha\beta u}_{\phantom{\dagger}\gamma}$ 
represent application of the anyonic fusion rules, and may be associated with vertices of the splitting and fusion trees. Under the isotopy invariance convention there is an additional factor of $[d_{c}/(d_{a}d_{b})]^\frac{1}{4}$ associated with the fusion of %vertex combining 
charges $a$ and $b$ into $c$, where $d_{x}$ is the quantum dimension of charge $x$, and similarly for splitting, but we will account for these factors separately. Thus constructed, the tensors $\tilde N$ satisfy $\tilde N^\gamma_{\alpha\beta u}\tilde N^{\dagger\alpha\beta u}_{\p\dagger\epsilon}=\delta_\epsilon^{\p\epsilon\gamma}$.

% when $N^c_{ab}\leq 1~\forall~\{a,b,c\}$. When some $N^c_{ab}>1$, $\tilde N^\gamma_{\alpha\beta}\tilde N^{\dagger\alpha\beta}_{\p\dagger\epsilon}$ will %also 
%contain a number of $n\times n$ blocks whose entries all take the value $n^{-1}$. 

%As these blocks act on homogeneous subspaces consisting of $n$ indistinguishable copies of the same charge, their actions are equivalent to the identity. We note that a more elegant construction is possible by defining $N'^c_{\p\prime ab}$ such that $N'^c_{\p\prime ab}=1~\forall~N^c_{ab}\geq 1$, with all other entries zero. If this is used throughout the preceding material in lieu of $N^c_{ab}$, then $\tilde N$ once again satisfies $\tilde N^\gamma_{\alpha\beta}\tilde N^{\dagger\alpha\beta}_{\p\dagger\epsilon}=\delta_\epsilon^{\p\epsilon\gamma}$. It should be noted that under this latter construction, the factor of $[d_{c}/(d_{a}d_{b})]^\frac{1}{4}$ associated with a fusion or splitting vertex is replaced by a factor of $[d_{c}/(d_{a}d_{b})]^\frac{1}{4}(N^c_{ab})^{-\frac{1}{2}}$.

When used as a representation of the fusion rules, the generalised multiplicity tensor $\tilde N^\gamma_{\alpha\beta u}$ and its conjugate $\tilde N^{\dagger\alpha\beta u}_{\phantom{\dagger}\gamma}$ permit us to increase or decrease the number of multi-indices on a tensor in a manner which is consistent with the fusion rules of the quantum symmetry group. This process is reversible provided the symmetry group is abelian or, for a nonabelian symmetry group, provided the total number of multi-indices on the tensor does not at any time exceed three.
%\footnote{We are assuming that when}
In constructing and manipulating a tensor network for a system of anyons, we will require only objects which respect the fusion rules of the anyon model. It is a defining property of such objects that when the number of multi-indices they possess is reduced to 1 by repeated application of $\tilde N$ and $\tilde N^\dagger$, non-zero entries may be found only in the vacuum sector. We imposed this requirement for states in \sref{sec:anyonstates}, and it is equivalent to the restriction we imposed on anyonic operators in \sref{sec:anyonops}.
In \rcite{singh2009} an equivalent condition was observed for tensors remaining unchanged under the action of a Lie group, and these tensors were termed \emph{invariant}. When working with invariant tensors, we may separately evaluate the components of the tensors acting on the degeneracy spaces (e.g. the nonzero blocks of $M_\alpha^{\phantom{\alpha}\beta}$), and the factors arising from loops and vertices of the associated spin network. This property greatly simplifies the contraction of pairs of tensors. 

In addition to increasing or decreasing the number of legs of a tensor, we may also use $\tilde N$ to ``raise'' the matrix representation of an operator from the space of $x$ sites to the space of $(x+x')$ sites. This is shown in \fref{fig:raiseoperator}, and the matrix representation of the raised operator is given by
\begin{equation}
M_{\p{\prime}\alpha}^{\prime\p{\alpha}\beta} = M_\gamma^{\p{\alpha}\delta}\tilde N^{\dagger \gamma\epsilon u}_{\p{\dagger}\alpha} \tilde N^\beta_{\delta\epsilon u}\label{eq:Mprime}
\end{equation}
where multi-index $\epsilon$ describes the fusion space of all sites in $(x+x')$ but not in $x$. 
Because the numeric factors associated with loops, vertices, and braiding (where applicable) are handled separately,
no factors of quantum dimensions appear in Eq.~\eref{eq:Mprime}.
\begin{figure}
\includegraphics[width=246.0pt]{raiseoperator}
\caption{``Raising'' of an operator $\hat M$ from sites $x$ to sites $x+x'$: (i) Operator $\hat M$ defined only on sites denoted $x$. (ii) Resolutions of the identity are inserted above and below $\hat M$, being constructed from tensors $\tilde N$ and $\tilde N^\dagger$. The central portion of this diagram is identified as corresponding to the new matrix $M'^{~\beta}_\alpha$ which describes $\hat M$ on $x+x'$. (iii) Loop and vertex factors in the central region are evaluated separately and eliminated. (iv) The tensor network corresponding to the new central portion is contracted. The $\tilde N$ and $\tilde N^\dagger$ tensors outside the central region become vertices of the fusion and splitting trees associated with $M'^{~\beta}_\alpha$. Together the trees and the matrix $M'^{~\beta}_\alpha$ constitute the raised version of $\hat M$. \label{fig:raiseoperator}}
\end{figure}

To act an operator $\hat M$ on a state $|\psi\ra$ in the matrix representation, we therefore connect the diagrams for $\hat M$ and $|\psi\ra$, eliminate all loops, and then raise the  matrix representation of the operator $\hat M$ using Eq.~\eref{eq:Mprime}, repeatedly if necessary, until the resulting matrix $M_{\p{\prime}\alpha}^{\prime\p{\alpha}\beta}$ may be applied directly to the state vector $c^\alpha$. 
Similarly it is possible to combine the matrix representations of operators, by connecting their diagrams appropriately, eliminating loops, and performing any required raising so that both operators act on the same fusion space. Their matrix representations can then be combined to yield the matrix representation of the new operator:
\begin{equation}
M^{(1\times2)\p{\alpha}\beta}_{\p{(1\times2)}\alpha} = M^{(1)\p{\alpha}\gamma}_{\p{(1)}\alpha} M^{(2)\p{\gamma}\beta}_{\p{(2)}\gamma},\label{eq:contractmatrices}
\end{equation}
and the fusion/splitting tree associated with this new operator is obtained as shown in \fref{fig:raiseoperator}.

Note that as yet, we have not described how two objects may be combined if their multi-indices are both up or both down, and are connected by a curved line. To contract such objects together, it is necessary to understand how bends act on the central matrix of an operator. Once this is understood, the bend can be absorbed into one of the central matrices, so that the connection is once again between an upper and a lower multi-index as in Eq.~\eref{eq:contractmatrices}. This process is described in \sref{sec:anyonmanip}.

\subsection{Manipulation of anyonic operators\label{sec:anyonmanip}}

As observed in \sref{sec:siteexpanyop}, when we describe a system entirely in terms of objects invariant under the action of the symmetry group, we may account separately for the numerical normalisation factors associated with the spin network. %, and the invariant tensors whose blocks inhabit the degeneracy spaces. 
However, as well as affecting these numerical factors, transformations of the fusion or splitting tree of an anyonic operator will typically also generate unitary matrices which act on the matrix representation of the operator. These matrices respect the symmetry of the anyon model, and thus can be written as block-diagonal matrices where each block is %individually 
a unitary matrix acting on %the degeneracy space associated with 
a particular charge sector. %, and may be represented diagramatically by the insertion of a grey circle on the appropriate leg of the . 
In terms of the diagrammatic notation of \sref{sec:tensordiagrammatic}, $F$ moves and braids therefore result in the insertion of a unitary matrix, as shown in \fref{fig:anyonopmanip1}. These matrices, whose entries are derived from the tensors $(F^{abc}_{d})_{(euv)(fu'v')}$ and $R^{ab}_c$ respectively, are raised if required, as described in \sref{sec:siteexpanyop}, and then contracted with $\Mop$, the matrix representation of the operator. 
\begin{figure}
\includegraphics[width=246.0pt]{newFmoveetc}
\caption{(i) $F$ move, and (ii) braiding, performed on a section of fusion tree in the diagrammatic notation of \protect{\sref{sec:tensordiagrammatic}}.\label{fig:anyonopmanip1}}
\end{figure}%
To compute the unitary matrices involved, it suffices to recognise that $F$ moves and braids are unitary transformations in the space of labelled tree diagrams.
Identifying the leg on which the unitary matrix is to be inserted, 
%Then, from the vertices involved in the manipulation, select the one which is closest to the central object of the operator. If this vertex is understood as a vertex tensor $\tilde N$, 
%, and that %the charge and degeneracy indices on the 
%when the vertex which is closest to the central object of the operator is understood as a tensor $\tilde N$,
%then 
the relevant region of the space of labelled diagrams is then enumerated by the multi-index which can be associated with this leg (compare \fref{fig:exampleedges}).
%of the %indicated 
%most proximal leg of $\tilde N$. % on which the unitary matrix is to be inserted. 
%Similar considerations apply on splitting trees, as is required by consistency with hermitian conjugation.
%The evaluation of a manipulation of the splitting tree is equivalent. %, provided the orderings of degeneracy indices on fusion and splitting spaces are consistent under hermitian conjugation.

Braiding is of particular importance when working in two dimensions, as an operator will necessarily be defined with respect to some arbitrary linear ordering of its legs, and when manipulating a tensor network it may be necessary to map between this original definition and other equivalent definitions, corresponding to different leg orderings. % as it is frequently associated with the process of linearising the 2-D lattice. %This is shown in the following example:
%The following example is particularly relevant to 2-D tensor networks: 
For example, let $\hat M$ be a four-site anyonic operator as shown in \fref{fig:absorbbraid}(i), which we wish to apply to a 2-D lattice. For the indicated linearisation of this lattice, application of $\hat M$ will require braiding as shown in \fref{fig:absorbbraid}(ii).
\begin{figure}
\includegraphics[width=246.0pt]{absorbbraid}
% Draw (i) in {1-D} and 2-D, with linearisation shown. All others in {1-D}.
\caption{(i) An operator $\hat M$ acting on sites on a 2-D lattice is defined with respect to some arbitrary linear ordering of these sites. (ii) When manipulating the tensor network, it may on occasion be computationally convenient for the lattice to be linearised according to some alternative linearisation scheme. In this example, the imposed linearisation scheme is indicated by the dotted line. To apply $\hat M$ to a different linearisation of the lattice may require braiding. The orientation of the braids can be determined by putting the fusion tree of (i) onto the 2-D lattice, then smoothly deforming the lattice into a chain in accordance with the linearisation prescription. (iii) The unitary matrices corresponding to the required $F$ moves and braiding operations may be absorbed into $\hat M$, defining a new operator $\hat M'$ on the linearised lattice.\label{fig:absorbbraid}}
\end{figure}%
By evaluating the unitary transformations corresponding to these braids and absorbing them into $\Mop$, we may define a new operator $\hat M'$ which acts directly on the linearised lattice without any intervening manipulations of the fusion/splitting trees.

We will also frequently wish to deal with tensor legs which bend vertically through 180 degrees. 
If working with an anyon model that has non-trivial Frobenius--Schur indicators, then 
indicator flags must be applied to all bends. %, and the initial assignment of these flags will be discussed in \sref{sec:constensnet}. This initial assignment is not particularly important as 
Like $F$ moves and braiding, the reversal of a Frobenius--Schur indicator flag is a unitary transformation%s may be reversed% during manipulation of tensor networks
, and once again this leads to the introduction of a unitary matrix which can be absorbed into a nearby existing tensor. %To complete our understanding of fusion and splitting trees, it is necessary only 
%In addition to this, 
%However, 
%to do this
%we may first need
However, we may wish to perform other operations on bends, such as absorbing them into fusion vertices or the central matrices of anyonic operators. We may also need to move a matrix $\Mop$ across a bend. We must therefore develop the description of bends in the new diagrammatic formalism.
%As these matrices will typically arise in the vicinity of bends, 
%to understand how bends act on fusion vertices and on the central matrices of anyonic operators.

In \rcite{bonderson2007} a prescription for absorbing bends into fusion vertices is given in terms of tensors $(A^{ab}_c)_{uv}$ and $(B^{ab}_c)_{uv}$, derived from the $F$ moves, and corresponding to clockwise and counter-clockwise bends respectively. 
%Although these tensors are derived from $F$ moves, %(see Appendix), 
%%it is more convenient to assign them a specific interpretation as described here. 
%%it is more convenient to provide a separate treatment of their action.
%%describe their action in the manner presented here.
%%their action admits a particularly concise treatment in the multi-index notation.
%our notation will benefit from providing them with an independent treatment.
The absorption of a clockwise or counterclockwise bend into a fusion vertex is reproduced in 
\fref{fig:anyonopmanip2}(i), and results in a vertex fusing upward- and downward-going legs.  
\begin{figure}
\includegraphics[width=246.0pt]{bending}
\caption{Vertical bending of legs (i) in the standard diagrammatic notation, and (ii) in the diagrammatic notation of \protect{\sref{sec:tensordiagrammatic}}. White triangles represent Frobenius--Schur indicator flags. (iii) Legs on the matrix representations of states and operators may also absorb bends. %(iv) 
\label{fig:anyonopmanip2}}
\end{figure}%
% shows, in the standard notation, how a clockwise or counterclockwise bend may be absorbed into a fusion vertex, resulting in a new type of vertex fusing an upward-going leg with a downward-going leg. 
We now assign new tensors $(\tilde N^\mrm{CW})^{\dagger\alpha u}_{\p\dagger\gamma\beta}$ and $(\tilde N^\mrm{CCW})^{\dagger\beta u}_{\p\dagger\alpha\gamma}$ to such vertices, such that writing these transformations in the notation of \sref{sec:tensordiagrammatic} is trivial. This is shown in \fref{fig:anyonopmanip2}(ii).

Explicit expressions for the new vertex tensors $(\tilde N^\mrm{CW})^\dagger$ and $(\tilde N^\mrm{CCW})^\dagger$ may be obtained by recognising that \fref{fig:anyonopmanip2}(i) describes the action of unitary transformations %$U_\delta^{\p\delta\gamma}$
on $\tilde N^{\dagger\alpha\beta u}_{\p\dagger\gamma}$. When the bend is counterclockwise, the corresponding unitary matrix is derived from $(A^{ab}_c)_{uv}$, and when the bend is clockwise, the unitary matrix is derived from $(B^{ab}_c)_{uv}$. We will denote these unitary matrices $A_\gamma^{\p\gamma\delta}$ and $B_\gamma^{\p\gamma\delta}$ respectively.
%
%, where $U$ is derived from $(A^{ab}_c)_{uv}$ when the bend is clockwise, or $(B^{ab}_c)_{uv}$ when the bend is counter-clockwise. 
%For the clockwise example, where $U_\delta^{\p\delta\gamma}$ is derived from $(A^{ab}_c)_{uv}$, we have
We then have
\begin{eqnarray}
(\tilde N^\mrm{CW})^{\dagger\alpha u}_{\p\dagger\gamma\beta}&=&A_\gamma^{\p\gamma\delta}\tilde N^{\dagger\alpha\epsilon u}_{\p\dagger\delta}\delta_{\epsilon\beta}\label{eq:Nbend}\\
(\tilde N^\mrm{CCW})^{\dagger\beta u}_{\p\dagger\alpha\gamma}&=&B_\gamma^{\p\gamma\delta}\tilde N^{\dagger\epsilon\beta u}_{\p\dagger\delta}\delta_{\epsilon\alpha}.\label{eq:Nbend2}
\end{eqnarray}
and conjugation describes equivalent vertices $\tilde N^{CW}$ and $\tilde N^{CCW}$ when a bend is absorbed into a splitting tree. 
%, noting that conjugation reverses the direction of the Frobenius--Schur indicators in \fref{fig:anyonopmanip2}.

Knowing how the absorption of bends acts on a vertex tensor, we may readily infer how the same process acts on the matrix representation of an operator. In \fref{fig:anyonopmanip2}(iii) we see a bend absorbed into the matrix $\Mop$, resulting in a new object with two lower multi-indices, $M'_{\alpha\beta}$. First we exploit the freedom to introduce fusion with the trivial charge (denoted $\mbb{I}$), with degeneracy 1. The corresponding $\tilde N^\dagger$ object takes only one value on its upper left multi-index, and is fully defined by $\tilde N^{\dagger\mbb{I}\beta 1}_{\p\dagger\gamma} = \delta^{\p\gamma\beta}_\gamma$. Absorbing the bend into this fusion vertex as per Eq.~\eref{eq:Nbend} yields $A_\gamma^{\p\gamma\delta}\delta_\delta^{\p\delta\epsilon}\delta_{\epsilon\beta} = A_{\gamma\beta}$, which may be then combined with $\Mop$ to give
\begin{equation}
M'_{\alpha\beta} = M_\alpha^{\p\alpha\gamma}A_{\gamma\beta}.
\end{equation}
In conjunction with the relationships given in \fref{fig:frobeniusschur}, this gives us the ability to move a matrix past a bend. 
\begin{figure}
\includegraphics[width=246.0pt]{frobeniusschur}
\caption{Opposing pairs of Frobenius--Schur indicators (i) on a pair of bends equivalent to the identity, and (ii) on a pair of bends such as might be used when computing a quantum trace. (iii) As a% gauge can always be chosen where
n anyon model can always be specified such that the Frobenius--Schur indicators are $\pm 1$, reversing a pair of contiguous opposed Frobenius--Schur indicator flags is always free%, for example going from (i) to (iii)
.
%Recommended initial choices of Frobenius--Schur indicators. (i) To allow a pair of bends to mutually annihilate, the Frobenius--Schur indicator flags should be opposed. (ii) For quantum traces and related processes, both Frobenius--Schur indicator flags should again be opposed, pointing in the same direction (either left or right).
\label{fig:frobeniusschur}}
\end{figure}%
An example of this is given in \fref{fig:movepastbend}, for which $M$ and $M'$ are related according to
\begin{equation}
M'^{\p\alpha\beta}_\alpha = A_{\alpha\gamma}M_\delta^{\p\delta\gamma}\varkappa_\epsilon^{\p\epsilon\delta}B^{\dagger\epsilon\beta}\label{eq:movepastbend}
\end{equation}
where $\varkappa_\epsilon^{\p\epsilon\delta}$ represents reversal of the Frobenius--Schur indicator flag on the lower bend.
\begin{figure}
\includegraphics[width=246.0pt]{movepastbend}
\caption{Moving a matrix across a bend in a tensor network diagram. (i) Initial diagram. (ii) Bends are absorbed into the matrix. (iii) New bends are introduced, in accordance with \protect{\fref{fig:frobeniusschur}(i)}. (iv) A pair of contiguous, opposed Frobenius--Schur indicators are reversed, as per \protect{\fref{fig:frobeniusschur}(iii)}. The initial and final matrices $M$ and $M'$ are related as specified in Eq.~\protect{\eref{eq:movepastbend}}.\label{fig:movepastbend}}
\end{figure}%
%where we have also used the relationship given in \fref{fig:frobeniusschur}(i). %, relationships of this sort enable the movement of matrices across bends when bringing together two operators, or when absorbing into an operator a unitary matrix which has been generated by a manipulation on a fusion or splitting tree.
Finally, bending may also allow more efficient contraction of pairs of anyonic operators, as shown in \fref{fig:controps_bends}.
\begin{figure}
%\begin{center}
\includegraphics[width=246.0pt]{controps_bends}
\caption{The use of bends may permit the more efficient contraction of pairs of anyonic operators. In the sequence of events marked (i), operator $\hat A$ is first raised to the space of three sites then contracted with $\hat B$. In sequence (ii) the operators are instead contracted using bends. For many anyon models the latter approach offers a significant computational advantage.\label{fig:controps_bends}}
%\end{center}
\end{figure}%


Having described the action of bends, it is customary also to introduce a second type of $F$ move which is described by the tensor $(F^{a_1a_2}_{a_3a_4})_{(a_5u_1u_2)(a_6u_3u_4)}$ (\fref{fig:newfmove}). 
\begin{figure}
%\begin{center}
\includegraphics[width=246.0pt]{newfmove}
\caption{Now that we may bend legs up and down it is customary to introduce a further type of $F$ move, derived by applying bends to the one presented in \protect{\fref{fig:anyons}}(iii). %Its action is presented here for completeness.
\label{fig:newfmove}}
%\end{center}
\end{figure}%
This tensor may be derived from $(F^{a_1a_2a_3}_{a_4})_{(a_5u_1u_2)(a_6u_3u_4)}$ by bending, and as with $(F^{a_1a_2a_3}_{a_4})_{(a_5u_1u_2)(a_6u_3u_4)}$ these $F$ moves perform a transformation of the fusion tree, accompanied by the introduction of a unitary matrix which can be absorbed into the matrix representation of the operator. These unitary matrices %are trivially seen to 
correspond to the consecutive application of a bend, an $F$ move of the original type% given in \fref{fig:anyons}(iii)
, and a second bend whose action is the inverse of the first.



% at the level of the degeneracy indices. 
%
%will in general affect not only these normalisation factors, but will affect also the degeneracy spaces of the operator matrix with which the fusion tree is associated. 
%
%In this section apply equally well to splitting trees.
%
% at the root of the fusion tree being transformed. Identical considerations apply for splitting trees, and although in this section we will provide illustrative examples for fusion trees only, for splitting trees the equivalent considerations apply.
%
%
%
%For anyonic systems, these operations on the spin network correspond to operations on the splitting or fusion trees, such as braiding, bending, and other deformations consistent with diagrammatic isotopy. These operations are well-understood and we direct the reader to the literature for details \cite[e.g.][]{kitaev2006,bonderson2007}. 
%
%However, under many circumstances it may be preferable to 


%We mention them here only to note that when a spin network manipulation acts entirely on the legs of some operator $\hat M$, the manipulation of the fusion/splitting tree may, if desired, be absorbed into $\hat M$ to define a new operator, $\hat M'$. 

%\subsection{Anyonic superoperators}
%
%Having developed the formalism of bends, we are now in a position 





% Plan: Role of bends in contracting operators.
% i) Superoperators <-> operators; 3-legged ws
% ii) Too-thin middle legs; two square ops, and 2-legged ws
% iii) More efficient contractions

%\subsection{Contraction of anyonic operators\label{sec:controp}}
%
%In \sref{sec:siteexpanyop} we saw how an anyonic operator can be applied to a state by connecting their diagrams, performing manipulations on the fusing and splitting trees, raising the matrix corresponding to the operator, and finally evaluating $c^{\prime\alpha}=c^\beta M_\beta^{\p\beta\alpha}$. Similarly, it is possible to contract together two anyonic operators. A simple case is shown in \fref{fig:controp}(i). 
%\begin{figure*}
%%\includegraphics[width=246.0pt]{}
%\caption{Contraction of pairs of anyonic operators. In these diagrams, all lines are assumed to correspond to indices of equal dimension. (i) $\hat A$ is a two-site operator and $\hat B$ is a 3-site operator. They may be contracted by raising $\hat A$ to three sites, eliminating the trees which lie between the matrix representations of these operators, then contracting these matrices as per Eq.~\protect{\eref{eq:controps}}. (ii) Alternatively, one leg of $\hat B$ may be bent down prior to contraction for a possible improvement in computational cost. (iii) If $\hat A$ and $\hat B$ are both two-site operators, acting as shown, then we could choose to first raise $\hat B$ to the space of three sites, denoted $\hat{\mbb{I}B}$, before contracting with $\hat A$ as in (ii). Operator $\hat A$ acts on the fusion space of two sites, indicated by $*$. (iv) Similarly, we could first raise $\hat A$ to the space of three sites, and then $\hat B$ acts on the two-site splitting space denoted $\dagger$. (v) Performing both bends is prohibited. The second bend attempts to factorise the fusion space acted on by $\hat A$, which is generally not possible for nonabelian anyon models. Diagrams~(vi) and (vii) demonstrate the application of these techniques to more complicated situations. In (vi) a superoperator is constructed from two operators, and in (vii) this superoperator is applied to an operator. Both diagrams use bends so that treatment of the superoperator is essentially no different to the contraction of operators which we have already described.\label{fig:controp}}
%% Show internal structure in diagrams.
%% (i): B 3 sites, A 2 sites & above. Requires an F move.
%% (ii): Same setup, with bend.
%% (iii): A(I) above, 2 sites. (I)B below, 2 sites. Raise B then do as (ii).
%% (iv): Same, but raise A then do as (ii).
%% (v): Prohibited double-bend. Do (iii) bend first, then put a X on arrow for doing (iv) bend.
%% (vi): Contraction of 2-site w and w^\dagger to give a 1-site superoperator.
%% (vii): Application of the superoperator to an operator.
%\end{figure*}%
%The fusion tree of operator $\hat A$ acts on a subspace of the splitting tree of operator $\hat B$, and the two operators can be combined by performing an $F$ move, eliminating a loop, raising the matrix $B$ to act on the fusion space of three sites rather than two, and then writing 
%\begin{equation}
%A_{\p\prime\alpha}^{\prime\p\alpha\beta}=B_\alpha^{\p\alpha\gamma}A_{\gamma}^{\p\gamma\beta}. \label{eq:controps}
%\end{equation}
%This process is directly analogous to acting an operator on a state. Alternatively, we may use bends to reduce the computational cost as shown in \fref{fig:controp}(ii). For a spin system, the cost of this contraction scheme would be less than that of approach (i), and the same also holds true for many anyon models.
%
%Consider now \fref{fig:controp}(iii). To combine operators $\hat A$ and $\hat B$, we could raise $\hat B$ to the fusion space of three sites, the raised operator here being denoted $\hat{\mbb{I}B}$. Operator $\hat A$ then acts on the space described by the fusion of two sites, indicated by $*$. We can then use bending to efficiently contract this network as in the previous example. Equivalently, as in (iv), we could raise operator $\hat A$ to the fusion space of three sites, then bend one of its legs before merging in $\hat B$. What we cannot choose to do, however, is to perform both bends before the contraction, as shown in \fref{fig:controp}(v). After the first bend has been performed, an implicit choice has been made: Either operator $\hat A$ is understood to be acting on the fusion space indicated by $*$, or $\hat B$ is understood to be acting on the splitting space indicated by $\dagger$. The second bend then factorises this space. For specificity, let us assume that the first bend has been performed as in (iii). Operator $\hat A$ now acts on the fusion space of two sites, $*$. If we denote the dimension of the fusion space of $n$ sites by $d_n$, then fusion space $*$ is of dimension $d_2$, which satisfies
%%this fusion space by $d_2$ and the dimension of a single site by $d_1$, then in general we have 
%$d_2\geq (d_1)^2$, where $d_2>(d_1)^2$ is possible for nonabelian anyon models. 
%The second bend divides the two-site fusion space on which $\hat A$ is to act into two one-site fusion spaces, and constitutes an effective projection of $A_\alpha^{\p\alpha\beta}$ onto a subspace %rank 
%of dimension $(d_1)^2$. For nonabelian anyons, this will typically yield an incorrect result. 
%%By performing the second bend, we divide the fusion space on which $\hat A$ is to act into two subspaces of dimension $d_1$. This effectively projects $A_\alpha^{\p\alpha\beta}$ onto a subspace of rank $(d_1)^2$, which for nonabelian anyons will generally yield an incorrect result. 
%
%%Consider now \fref{fig:controp}(iii). To combine operators $\hat A$ and $\hat B$, we could raise $\hat B$ to the fusion space of three sites, the raised operator here being denoted $\hat{\mbb{I}B}$. Operator $\hat A$ then acts on the space described by the fusion of two sites, indicated by $*$. We can then use bending to efficiently contract this network as in the previous example. Equivalently, as in (iv), we could raise operator $\hat A$ to the fusion space of three sites, then bend one of its legs before merging in $\hat B$. What we cannot choose to do, however, is to perform both bends before the contraction, as shown in \fref{fig:controp}(v). After the first bend has been performed, an implicit choice has been made: Either operator $\hat A$ is understood to be acting on the fusion space indicated by $*$, or $\hat B$ is understood to be acting on the fusion space indicated by $\dagger$. %The second bend then requires that this fusion space. 
%%For specificity, let us assume that the first bend has been performed as in (iii). Operator $\hat A$ now acts on the fusion space of two sites. If we denote the dimension of the fusion space of $n$ sites by $d_n$, then fusion space $*$ is of dimension $d_2$, satisfying
%%%this fusion space by $d_2$ and the dimension of a single site by $d_1$, then in general we have 
%%%We have
%%$d_2\geq (d_1)^2$. Having $d_2>(d_1)^2$ is possible for nonabelian anyon models. 
%%
%%To perform the second bend, and move the affected leg from the lower (splitting) tree of $\hat A$ to the upper (fusion) tree, it is necessary to separate the splitting tree of $\hat A$ into two separate legs, one of which will absorb the bend. %This is done by absorbing a vertex (represented by $\tilde N^\dagger$) into $A\alphabeta$. 
%%The entire bending process is reversible on $\hat A$, as the number of legs does not exceed three, but has the additional effect of also dividing the fusion space $*$ into two subspaces each of dimension $d_1$. %For nonabelian anyons, p
%%Performing such a splitting prior to contracting operator $\hat A$ with $\hat B$
%%%
%%%, corresponding also to the factorisation of the fusion space $*$. 
%%%
%%%The bending process shown is reversible on $\hat A$, but in the process of moving this leg from the lower (splitting) tree of $\hat A$ to the upper (fusion) tree, 
%%%
%%%
%%%Moving one of the legs from the splitting tree of $\hat A$ to the fusion tree of $\hat A$ 
%%%%Dividing the fusion space on which $\hat A$ is to act 
%%%implies that the splitting tree may be safely divided, 
%%constitutes an effective projection of $A_\alpha^{\p\alpha\beta}$ onto a rank $(d_1)^2$ subspace, 
%%which will for nonabelian anyons generally yield an incorrect result.
%%%yielding an incorrect result. 
%
%To avoid the accidental introduction of such projections, it suffices to ensure that when multiplying together the matrix representations of the two operators as in Eq.~\eref{eq:controps}, for one of the free multi-indices $\alpha$ and $\beta$ the degeneracies of every charge in that index are less than or equal to the degeneracy of the same charge in $\gamma$.
%
%Provided this rule is adhered to, it is possible to use bends to efficiently evaluate even relatively complicated contractions such as those shown in \fref{fig:controp}(vi) and (vii). %, which corresponds to the construction and application of the scaling superoperator of the ternary 1-D MERA.
%
%%In general it may therefore be impossible to 
%%
%% and for nonabelian anyons it may be impossible to factorise 
%%
%%, and recover the previous example, subsequently contracting as per (ii). Alternatively, we could raise operator  
%%
%
%
%
%%
%%\subsection{Creating anyonic superoperators\label{sec:superop}}
%%
%%There is one further special situation to consider, which is the contraction of two operators to yield a superoperator, such as the structure shown in \fref{fig:superoperator}(i).
%%\begin{figure}
%%\begin{center}
%%%\includegraphics[width=246.0pt]{}
%%\caption{(i) Original tensor network. (ii) \& (iii) Two possible structures of the constituent anyonic tensors. (iv) \label{fig:superoperator}}
%%% @@@ Note all legs have same charges and degeneracies
%%\end{center}
%%\end{figure}%
%%In this situation the fusion and splitting trees of the operators may not naturally describe a raising process for the operator matrices. Alternatively, the anyonic structure of the tensors may be such that the contraction appears trivial, except that the degeneracies of the summed index are smaller than those of either free index in Eq.~\eref{eq:contractmatrices}. The resolution of these problems are presented in the figure, exploiting a mixture of bending and vertex tensors. Also note that bending may be used to decrease the cost of contracting a tensor network: Following diagram (@@@), an alternative approach would have been to raise matrix @@@ to have three lower legs and two upper legs, before contracting with $\mbb{I}$. However, the computational cost would have been higher than that of the approach shown.
%%
%%These approaches are not restricted to the evaluation of superoperator diagrams. For example, consider \fref{fig:smalllink}.
%%\begin{figure}
%%\begin{center}
%%%\includegraphics[width=246.0pt]{}
%%\caption{\label{fig:smalllink}}
%%% @@@ Note all legs have same charges and degeneracies
%%\end{center}
%%\end{figure}%
%%This figure shows a simple network requiring some care to contract. Attempting to use bends as in diagram (ii) results in the degeneracies of $\gamma$ in Eq.~\eref{eq:contractmatrices} being smaller than those of $\alpha$ and $\beta$. This problem may be circumvented either by first raising one of the tensors to have three upper and three lower legs (diagrams (iii)-(@@)), or by identifying regions of the diagram with bending and identity tensors as shown in diagrams (@@)-(@@). Careful application of all these techniques can be important in determining the optimal contraction of a tensor network.
%%
%%%
%%%first exploiting diagrammatic isotopy, and then identifying portions of the tensor network diagram with tensors which may be constructed from the identity matrix and from bends. Some ingenuity may be required to find the optimal contraction technique for such diagrams. The same technique may be also used to circumvent the caveat attached to \eref{eq:contractmatrices}. Figure~\ref{fig:smalllink}(i) shows a simple diagram which may easily be contracted by first raising both of the tensors to the fusion space of three sites (diagrams (ii)-(iii)). If bends have been applied to these matrices (diagrams (iv)-(v)) the resolution is less obvious, but may be achieved by the technique of constructing simple tensors which correspond to portions of the fusion and splitting tree diagrams. 
%%%
%%%. Here, diagram (i) can be safely contracted by first raising either tensor onto the fusion space of three sites (the upper sequence of diagrams). 

\subsection{Constructing a tensor network\label{sec:constensnet}}

Now that we have developed a formalism for anyonic tensors,
%tensor network formalism for anyonic variables, 
we may convert an existing tensor network algorithm for use with anyons. 
First, the tensor network must be drawn in such a manner that every leg has a discernible vertical orientation. Although these orientations may be changed during manipulation of the tensor network, an initial assignment of upward or downward direction is required. Second, all tensors must be represented by entirely convex shapes, such as circles or regular polygons. For existing tensor network algorithms such as MERA and PEPS, this requirement is trivial. However, it is conceivable that future algorithms might involve superoperator-type objects whose graphical representations interleave upward- and downward-pointing legs. Concavities on these objects may be eliminated by replacing some of their upward-pointing legs with downward-pointing legs (or vice versa), followed by a bend (\fref{fig:generalTN}(i)-(ii)). 
\begin{figure}
\includegraphics[width=246.0pt]{generalTN}
\caption{Construction of an anyonic tensor corresponding to a normal tensor with more than three legs. (i) The original tensor. (ii) If required, any concavities are eliminated by introducing bends. (iii) Frobenius--Schur indicators are assigned to the bends. (iv) Directions are assigned to all legs, consistent with the rest of the network. %Where these legs represent shared indices, the direction must be assigned consistently on both of the tensors which share the index. 
(v) Legs are collected together into fusion and splitting trees. The central object, representing degrees of freedom of the tensor, now has less than four legs. (vi) If desired, bends can be re-absorbed into the fusion and splitting trees.
\label{fig:generalTN}}
\end{figure}
A similar treatment may be applied to any superoperators which arise during manipulations of the tensor network, introducing a pair of bends as in \fref{fig:frobeniusschur}(i) and then absorbing one into the matrix representation of the object.

If working with an anyon model that has non-trivial Frobenius--Schur indicators, then indicator flags must %now 
be applied to all bends. 
%When applying Frobenius--Schur indicators to bends, the 
Initial choices are a matter of convenience, and it is frequently possible to assign these indicators in opposed pairs, as shown in \fref{fig:frobeniusschur}. If these paired indicators are not flipped or are only flipped in adjacent opposed pairs %at all 
during subsequent manipulations of the tensor network, then they may %even
%often
frequently be left implicit.

Next, if there exist charges in the anyon model which are not self-dual, a direction (represented by a solid arrow) must be assigned to every multi-index. Any tensor with more than three legs (e.g. $M$ in \fref{fig:generalTN}) is then replaced by %rewritten as 
a trivalent tensor network
consisting of a core object, e.g. $M\alphabeta$, which contains the free parameters of the tensor, and as many copies of $\tilde N$ or $\tilde N^\dagger$ as are required to provide the correct output legs. 
These tensors $\tilde N$, $\tilde N^\dagger$ %then 
correspond to vertices in the fusion and splitting trees associated with $M\alphabeta$, yielding the corresponding anyonic tensor. Objects with three legs or less can be directly identified with an anyonic tensor object carrying the appropriate number of indices (i.e. three multi-indices and a vertex index $u$), though for consistency with the methods described in Sections~\ref{sec:siteexpanyop} and \ref{sec:anyonmanip} we point out that it is possible to similarly replace three-legged objects with anyonic operators consisting of a central matrix $\Mop$ and a fusion or splitting vertex, if desired.

Any bends introduced earlier may now be reabsorbed, so that some vertices now correspond to $(\tilde N^\mrm{CW})$, $(\tilde N^\mrm{CCW})$, $(\tilde N^\mrm{CW})^\dagger$, and $(\tilde N^\mrm{CCW})^\dagger$. This step, however, is optional as it may be more convenient for subsequent manipulations of the tensor network if the bends are left explicit. 
%If a tensor repeatedly interleaves upward- and downward-pointing legs% as in \fref{fig:generalTN}(iv)
%, it may be necessary to replace some upward-pointing legs with downward-pointing legs followed by bends, or vice versa, %before constructing its fusion and splitting trees, 
%in order to permit a construction in which the central object has three indices or less. Some ambiguity exists in the choice of which legs to bend on such an object, but this is unimportant as the different possibilities may be freely interconverted by means of the manipulations described in \sref{sec:anyonmanip}.
The anyonic tensors are 
%Once anyonic equivalents have been constructed for all tensors in the network, these
%The resulting anyonic tensors are 
then connected precisely as in the original ansatz. %Exceptionally, bend-derived vertex tensors such as $\tilde N^\mrm{CW}$ may also be required, although this is unusual in the initial construction and is necessary only if a tensor repeatedly interleaves upward- and downward-pointing legs, so that vertices of these types are required in order to keep the total number of indices on the central object to three or less.

%Where a leg curls through $360^\circ$, this may subsequently result in the performance of a quantum trace, and it is preferable that once again the indicator flags should be opposed (\fref{fig:frobeniusschur}(ii)). For all other bends, the choice is unimportant.

%
%as in \fref{fig:frobeniusschur}.
%Where diagrammatic isotopy indicates that a pair of vertical bends can annihilate one another, the flags should be opposed as in \fref{fig:frobeniusschur}(i), while
%all other vertical bends may be associated with Frobenius--Schur indicator flags pointing to the right. For a pair of bends having a combined curvature of $360^\circ$, for example in a quantum trace, this is the natural choice of orientation. 
%For the remaining single bends there is no natural choice of Frobenius--Schur indicator, and for definiteness, we choose that these indicators will also be directed to the right. If required, later reversal of a Frobenius--Schur indicator will simply generate a unitary transformation on the affected index which can always be absorbed into a nearby tensor.

% manipulated in an equivalent manner, 
Manipulations of the anyonic tensor network are equivalent to those performed on the spin version of the ansatz, differing only
%differing from the spin version of the ansatz only 
in that the degrees of freedom of the tensor network are now expressed entirely by the at-most-trivalent central objects, and certain topological elements such as braids and vertical bends must be accounted for in accordance with the prescriptions of \sref{sec:anyonmanip}. 
These changes may naturally imply minor changes to the manipulation algorithms, and we will see examples of this in the {1-D} MERA. Similar considerations will %also 
apply %for
to other tensor network algorithms. %As the unitary transformations associated with manipulations of the fusion and splitting trees are independent of the content of the central object of an anyonic tensor, these transformations need only be computed once, 



%%the variable content consists of only the 
%%at-most-trivalent central 
%%objects, %and so 
%%for which 
%%the
%%manipulation algorithms must be amended accordingly. 
%Some modifications may consequently have to be made to the manipulation algorithms, and %though these are typically straightforward,
%%The manipulation algorithms are amended accordingly, 
%we will see %an 
%examples of this %will be seen 
%for the 1-D MERA; %. %, where instead of computing the environment of the entire anyonic object $\hat{u}$ derived from tensor $u$, instead the fusion and splitting trees associated with the core object $u\alphabeta$ are also absorbed into the environment, yielding an environment not for $\hat{u}$, but for $u\alphabeta$, and hence the update procedure directly targets the optimisable content of $\hat{u}$. 
%%Similar considerations will apply to other tensor network algorithms.
%similar considerations will apply 
%to other tensor network algorithms. %However, the tools of \sref{sec:anyonmanip} 

%These modifications merely constitute the application of the techniques described in \sref{sec:anyonmanip} ($F$ moves, braiding, bends, and contraction with vertex tensors such as $\tilde N$) to implement the existing methods of the tensor network algorithm in a manner which is consistent with anyonic statistics.

Our construction of an anyonic tensor network draws upon two important elements which have previously been observed in other, simpler, physical systems:
\begin{enumerate}
\item Tensors in the ansatz exhibit a global symmetry, which may be nonabelian. Exploiting a nonabelian symmetry requires that the ansatz be written in the form of a trivalent tensor network. This has previously been observed and implemented for nonabelian Lie group symmetries such as $SU(2)$\citestop{singh2009}
\item Tensors in the ansatz must be able to account for non-trivial exchange statistics. This has previously been observed in the simulation of systems of fermions\citecomma{corboz2010,kraus2010,pineda2010,corboz2009,barthel2009,shi2009,pizorn2010,gu2010} where efficient implementation of particle statistics can be achieved through the use of ``swap gates''\citestop{pineda2010,corboz2009,barthel2009,pizorn2010} 
\end{enumerate}
In both cases, anyonic tensor networks extend the concepts introduced in previous work. The symmetry structure of an anyon model may be a quantum group, for example a member of the series $SU(2)_k$, $k\in\mbb{Z}^+$, rather than having to be a Lie group, %being restricted to Lie groups, 
and this permits representation of nonabelian anyonic systems whose Hilbert space does not admit decomposition into a tensor product of local Hilbert spaces. Similarly, anyonic braiding may be implemented using a generalisation of the fermionic ``swap gate'' formalism. When braiding, particle exchange may introduce transformation by a unitary matrix rather than by a sign, and efficient implementation of the resulting swap gates is particularly important for the simulation of 2-D systems. 

Although anyonic systems pose a number of unique challenges, we see that these are addressed by developments based on existing techniques, and we therefore anticipate that the resulting generalisations of existing tensor network ans\"atze %are 
should still be capable of accurately representing the states of an anyonic system.
% and of preserving (within the limits of the approximations inherent in the ansatz) anyonic degrees of freedom which admit no local representation.

%In addition, systems of nonabelian anyons


%tensor networks for systems of nonabelian anyons introduce a further feature which is accounted for by the structure 

%not previously observed, which is that the Hilbert space of the system does not admit decomposition into a product of local Hilbert spaces.

\subsection{Contraction of anyonic tensor networks}


The techniques described in Secs.~\ref{sec:siteexpanyop} and \ref{sec:anyonmanip} ($F$ moves, braids, bending of legs, elimination of loops, diagrammatic isotopy, flipping of Frobenius--Schur indicator flags, and the use of $\tilde N^{(\dagger)}$ tensors) suffice to contract any network of anyonic tensors written in the form of matrices with degeneracy indices, and unlabelled trees. Through careful application of these techniques, and avoiding at all times processes which would yield a tensor with more than three legs, the matrix representations of any pair of contiguous tensors in a network may always be brought into conjunction such that their multi-indices can be contracted in the manner of Eq.~\eref{eq:contractmatrices}, and any tensor network may be contracted by means of a sequence of such pairwise contractions.

That a tensor network may represent a system of anyons in this way
%This 
is possible because throughout the anyonic tensor network, each value of a degeneracy index is associated with a specific labelling of the corresponding unlabelled tree. Consequently it is always possible %at any time 
to fully reconstruct any operation in terms of the more verbose representation of \fref{fig:operators}.

An anyonic tensor network is therefore fully specified merely by the unlabelled tree (with Frobenius--Schur indicator flags if required), and the values and locations of the matrix representations of its tensors, written in the degeneracy index form. %Because it is always in principle possible to fully contract this network to specify a single object in the form of \fref{fig:blobtensors}, 


\section{Example: The {1-D} MERA}

\subsection{Construction\label{sec:MERAconstr}}

To construct an anyonic MERA for a {1-D} lattice with $n$ sites, where $n$ satisfies $n=2\times3^k,~k\in\mbb{Z^+}$, we begin with a ``top'' tensor on a two-site lattice $\mc{L}_\tau$ whose matrix representation is of a computationally convenient size. (The top tensor is named for its position in the usual diagrammatic representation of the MERA, where diagrams with open legs at the bottom correspond to a ket. For anyons the converse convention applies, and consequently in \fref{fig:MERAoperators}(i) the ``top'' tensor is ironically located at the bottom.)
%however, typically follow the converse convention where a tree with open legs at the top corresponds to a ket, and hence ironically, in \fref{fig:MERAoperators}(i), the top tensor may be found at the bottom.)
%Due to the customary orientation of diagrams representing anyonic states, however, in \fref{fig:MERAoperators}(i) it is, ironically, at the bottom.)
\begin{figure}
\includegraphics[width=246.0pt]{MERAoperators}
\caption{Construction of a {1-D} ternary MERA on a periodic lattice from anyonic operators. (i) The ``top'' tensor, $\hat{T}$. (ii) Isometries, $\hat{w}$. (iii) Disentanglers, $\hat{u}$. The fusion tree representing an anyonic state (or ket) is usually drawn with the lattice sites at the top, so this MERA has been constructed ``upside down'' when compared with the diagrams in Refs.~\protect{\onlinecite{evenbly2009,vidal2010}}. This is unimportant, and we could equally well have decided to follow the convention usually adopted in tensor network algorithms, labelled the tensors in (i)-(iii) by $T^\dagger$, $w^\dagger$, and $u^\dagger$, and identified diagram (i)-(iii) as a bra.
(iv) Structure of a 2-site term in the Hamiltonian, $\hat{h}$, or a 2-site reduced density matrix, $\hat{\rho}$. 
(v) Disentanglers and (vi) isometries satisfy the relationships $\hat u\hat u^\dagger=\mbb{I}$, $\hat w\hat w^\dagger=\mbb{I}$.
\label{fig:MERAoperators}}
\end{figure}

To each leg of the top tensor, we now append an isometry (\fref{fig:MERAoperators}(ii)). The matrix representations of the isometries consist of rectangular blocks, as described in \sref{sec:degenexp}, and we choose isometries whose fusion trees have three legs, so as to construct a ternary MERA\citestop{evenbly2009} Next, disentanglers are applied above the isometries.
For periodic boundary conditions this must be performed in a manner which respects the anyonic braiding rules, as shown in \fref{fig:MERAoperators}(iii). 
We identify the open legs of the resulting network as the sites of a lattice $\mc{L}_{\tau-1}$, and the rows of disentanglers and isometries may be understood as a coarse-graining transformation taking a finer-grained lattice $\mc{L}_{\tau-1}$ into a coarser-grained lattice $\mc{L}_\tau$, similar to the standard MERA. 
Note that the geometry of the periodic lattice is reflected by the connections of the disentanglers.
Specifically, whether the outside legs are braided over or under the other lattice sites reflects whether the lattice closes towards or away from the observer. 
 
The application of anyonic isometries and disentanglers is now repeated $k$ times (\fref{fig:MERAoperators}(i)-(iii) corresponds to $k=1$), until the ansatz has $n$ legs. The final row of isometries should be chosen such that each of their upper legs have the same charges and degeneracies as the sites of the physical lattice $\mc{L}_0$, and the open legs above the last row of disentanglers are identified with the physical lattice. For coarse-grained lattices $\mc{L}_1$ to $\mc{L}_\tau$, the dimensions of the lattice sites correspond to the lower legs of the isometries and are chosen for computational convenience, subject to the requirement that %the dimension of 
each charge sector is sufficiently large to adequately reproduce the physics of the low-energy portion of the Hilbert space. For all other legs, their charges and degeneracies are determined by requiring consistency with Eq.~\eref{eq:compounddegens}. Initial choices of which charges to represent on the ``top'' tensor and on the lower legs of the isometries, and with what degeneracies, must be guided either by prior knowledge about the physical system, or by balancing computational convenience against the inclusion of a broad and representative range of possible charges. When used in a numerical optimisation algorithm, the choice of relative weightings for the different charge sectors may often be refined by examination of the spectra of the reduced density matrices on the coarse-grained lattices, after initial optimisation of the tensor network is complete.

%it may be possible to further improve the result obtained by tuning the relative weighting of the different charge sectors, guided by the breadth of the spectra of the reduced density matrices on the coarse-grained lattices.

This concludes construction of the MERA for a state on a finite, periodic {1-D} anyonic lattice. That this tensor network does represent an anyonic state is easily seen by sequentially raising tensors, performing $F$ moves, and combining tensors, until the entire network is reduced to a single vector whose length is equal to the dimension of the physical Hilbert space, and an associated fusion tree. These then represent the state of the system as per Eq.~\eref{eq:statepsi}. The structure of this tensor network closely resembles that of the normal MERA, according to the identifications given in \fref{fig:MERAoperators}, and consequently we anticipate that it will share many of the same properties, including the ability to reproduce polynomially decaying correlators in strongly correlated physical systems.
Open lattices may also be easily represented by omitting the braided disentanglers at the edge of the diagram.

We also note that in common with the MERA for spins, the anyonic MERA may be understood as a quantum circuit, although one which carries anyonic charges in its wires. Any junction in the fusion/splitting trees may be associated with a $\tilde N$ or $\tilde N^\dagger$ tensor, and the entire network may be considered as the application of a series of gates to a Hilbert space of fixed dimension beginning mostly (or entirely, if the top tensor is considered to be the first gate) in the vacuum state, with individual gates introducing entanglement across some limited number of wires.


\subsection{Energy minimisation\label{sec:MERAoptimisation}}

The anyonic MERA can be used as a variational ansatz to compute the ground state of a local Hamiltonian. The Hamiltonian is introduced as a sum over nearest neighbour interactions, each term having the form of \fref{fig:MERAoperators}(iv), and 
optimisation of the tensor network is carried out in the usual manner\citestop{evenbly2009} Also as per usual, Hamiltonians involving larger interactions, such as next-to-nearest neighbour, can be accommodated by means of an initial exact $n$-into-one coarse-graining of the physical lattice.
%A Hamiltonian on the physical lattice is introduced in the form of \fref{fig:MERAoperators}(iv), with its matrix representation $H_\alpha^{\p{\alpha}\beta}$ having eigenvalues all less than or equal to zero,
%The Hamiltonian should be chosen so that the matrix representation of the 
%and regions of the anyonic operator tree are identified with tensors in a standard {1-D} MERA as shown in \fref{fig:MERAoperators}.

As in \rcite{evenbly2009}, optimisation of the MERA then consists of repeatedly lifting the Hamiltonian from $\mc{L}_0$ to the coarse-grained lattices, updating their isometries and disentanglers, and lowering the reduced density matrix, or the top tensor and its conjugate. When lifting the Hamiltonian or lowering the reduced density matrix, then the diagrams in \rcite{evenbly2009} taken in conjunction with the key given in \fref{fig:MERAoperators} serve to describe networks of anyonic operators which, when contracted to a single operator, yield the lifted form of the Hamiltonian or lowered form of the reduced density matrix respectively. Similarly, when optimising disentanglers or isometries, the diagrams of \rcite{evenbly2009} and the identifications in \fref{fig:MERAoperators} indicate how to construct an anyonic operator which constitutes the environment of the anyonic operator being optimised. However, once the admissible ranges of charges and degeneracies on each leg have been fixed, the only optimisable content of an anyonic operator is its matrix representation. 
Consequently, the fusion and splitting tree contributions should be evaluated and absorbed into the operator and its environment,
%also be absorbed into the environment operator.
%This reduces both the environment and the operator being updated 
reducing them both 
to their matrix representations, denoted $M$ and $E$ respectively (see \fref{fig:environment}).
\begin{figure}
\includegraphics[width=246.0pt]{environment}
\caption{(i) Anyonic operator $\hat E$ constitutes the environment of operator $\hat M$. %Frobenius--Schur indicator flags are . 
Factors arising from the fusion and splitting trees should be evaluated and absorbed into matrices $E$ and $M$, following which (ii) matrix $E$ constitutes the environment of matrix $M$. After (iii) updating the matrix $M$ %(with the new value denoted here by $u'$)
to $M'$, (iv) the fusion and splitting trees of $\hat M$ should be reinstated, the numerical factors associated with this process being the inverse of the fusion tree factors previously absorbed into matrix $M$. Frobenius--Schur flags in (i)-(ii) are represented by white triangles, and are not to be confused with the black arrows which indicate the orientation of lines in the fusion/splitting trees.\label{fig:environment}}
\end{figure}%
If the singular value decomposition of $E$ is written
$E=USW^\dagger$, then the updated matrix content $M$ of the anyonic operator being optimised is given by $-WU^\dagger$, minimising the value of $\mrm{Tr}(EM)$ subject to the usual constraint for disentanglers and isometries that $\hat M\hat M^\dagger=\mbb{I}$ %where $\hat M$ is a disentangler $\hat u\hat u^\dagger=\mbb{I}$ and $\hat w \hat w^\dagger=\mbb{I}$ 
(\fref{fig:MERAoperators}(v)-(vi)). 
%\footnote{Strictly, for quantum groups having multiplicities such that some $N^c_{ab}>1$, \protect{$\hat u\hat u^\dagger$} and \protect{$\hat w\hat w^\dagger$} are represented by block diagonal matrices consisting of blocks whose side lengths correspond to the values of $N^c_{ab}$. Where the side length of a block is $n$, the block constitutes an $n\times n$ matrix whose entries are all $n^{-1/2}$. However, all blocks having $n>1$ act on homogeneous subspaces corresponding to $n$ indistinguishable copies of the same charge, and thus their action is indistinguishable from the identity.}
The fusion/splitting tree content of the operator can then be restored, along with any appropriate numerical factors that may be required.

As with the standard MERA, the ``top'' tensor is constructed by diagonalising the total Hamiltonian on the most coarse-grained lattice, $\hat H_{\mrm{tot}}$ on $\mc{L}_\tau$. As $\mc{L}_\tau$ is a two-site lattice, the total Hamiltonian $\hat H_\mrm{tot}$ is a sum of two terms, $\hat H_{12}$ and $\hat H_{21}$. For the translation-invariant anyonic MERA, we may formally define $\hat H_{21}$ in terms of $\hat H_{12}$ as shown in \fref{fig:H21}, and the top tensor $\hat T$ (together with any factors arising from the chosen normalisation scheme) then corresponds to the lowest-energy eigenstate of $\hat H_\mrm{tot}$.
\begin{figure}
\includegraphics[width=246.0pt]{H21}
\caption{Definition of $\hat H_{21}$ in terms of $\hat H_{12}$, on the most coarse-grained lattice ($\mc{L}_\tau$) of the translation invariant periodic MERA. Lattice $\mc{L}_\tau$ is a two-site periodic lattice.\label{fig:H21}}
\end{figure}%

\subsection{Scale invariant MERA}

Having identified the anyonic counterparts of the tensors of the standard MERA, and described how these tensors may be lifted, lowered, and optimised, the algorithm for the scale-invariant MERA described in \rcite{pfeifer2009} may also be implemented for anyonic systems, simply by applying the dictionary of \fref{fig:MERAoperators} and the techniques described in \sref{sec:MERAoptimisation}. As with optimisation of $\hat{u}$ and $\hat{w}$, the computation of the top reduced density matrix (which is a descending eigenoperator of the scaling superoperator with eigenvalue 1) may be understood as a calculation of the matrix component $\rho_\alpha^{\p{\alpha}\beta}$ of the reduced density matrix $\hat{\rho}$. The ascending eigenoperators of the scaling superoperator, or local scaling operators of the theory, may also be computed in this manner.

\subsection{Results}

To demonstrate the effectiveness of the anyonic generalisation of the MERA, we applied it to a {1-D} critical system of anyons whose physical properties are already well known: The golden chain\citestop{feiguin2007} This model consists of a string of Fibonacci anyons subject to a local interaction. Fibonacci anyons have only two charges, $1$ (the vacuum) and $\tau$, and one non-trivial fusion rule ($\tau\times\tau\rightarrow 1+\tau$). %Both charges are self-dual, and both Frobenius--Schur indicators are 1.
%A full specification of the model may be found in \rcite{bonderson2007}.
The simplest local interactions for a chain of Fibonacci $\tau$ anyons are nearest neighbour interactions favouring fusion of pairs into either the $1$ channel (termed antiferromagnetic, or AFM), or the $\tau$ channel (termed ferromagnetic, or FM). Both choices correspond to critical Hamiltonians, associated with the conformal field theories $\mc{M}(4,3)$ and $\mc{M}(5,4)$ for AFM and FM couplings respectively. Individual lattice sites are each associated with a charge of $\tau$.
%It is worth remarking that the Fibonacci anyon model can also be understood as the even subspace of the anyon model derived from quantum group $SU(2)_3$. In turn, nearest neighbour interactions on the $SU(2)_k$ anyon models may be mapped to the 2-D classical Restricted Solid On Solid (RSOS) model with a transfer matrix invariant under the action of $SU(2)_k$ \cite{pasquier1990,sierra1997}. Examples of these RSOS Hamiltonians have been studied extensively using analytical techniques \cite{alcaraz1987}, variants of the DMRG algorithm \cite{sierra1997,tatsuaki2000}, and Quantum Monte Carlo \cite{todo2001}, but we will not exploit this special case mapping, and will instead continue to work explicitly in terms of Fibonacci anyons.

The AFM and FM Hamiltonians act on pairs of adjacent Fibonacci anyons. On a pair of lattice sites each carrying a charge of $\tau$, the matrix representations of the AFM and FM Hamiltonians are written
\begin{equation}
(H\alphabeta)_\mrm{AFM}=\left(\begin{array}{cc}\!\!-1\,\,&0\\\!\!\p{-}0\,\,&0\end{array}\right)
\quad
(H\alphabeta)_\mrm{FM}=\left(\begin{array}{cc}0&\p{-}0\\0&-1\end{array}\right)\label{eq:Hamiltonians}
\end{equation}
where a multi-index value of 1 corresponds to the vacuum charge, 2 corresponds to $\tau$, and the charges are non-degenerate. %We chose to analyse the antiferromagnetic case as detailed energy spectra for finite chains are available in the literature \cite{feiguin2007}.
%it is worth noting that the fusion rules for Fibonacci anyons correspond exactly to the even subalgebra of $SU(2)_3$, and these Hamiltonians may also be re-expressed as nearest neighbour interactions on a chain of spin-$\frac{1}{2}$ $SU(2)_3$ anyons \cite{}. There then exists a well-known mapping to the Restricted Solid On Solid (RSOS) model, 
%After optimising 
We optimised a scale-invariant MERA on the golden chain for each of these Hamiltonians% in turn
, %we
and computed local scaling operators using the tensor network given in \fref{fig:scalingsuperop}. 
\begin{figure}
\includegraphics[width=246.0pt]{scalingsuperop}
\caption{Determination of eigenoperators ($\phi$) and associated scaling dimensions ($\Delta$) for the one-site scaling superoperator of the anyonic {1-D} MERA. Eigenoperators may be classified according to the charges on edges $y_1$ and $y_2$. One interpretation of these labels is %as follows: In 
that, in
addition to the sites of the {1-D} lattice, there may exist free charges lying in front of and behind the anyon chain. The labels $y_1$ and $y_2$ then represent the transfer of charge between these regions and the {1-D} lattice. 
\label{fig:scalingsuperop}}
\end{figure}%
The operators calculated using this diagram may be classified according to the values of the charge labels $y_1$ and $y_2$, and the scaling dimensions and conformal spins which we obtained are given in Tables~\ref{tab:results} and \ref{tab:results2}, and \fref{fig:results}.
\begin{table}
\begin{tabular}{|c|c|c|}
\hline
\multicolumn{3}{|c|}{$y_1=y_2=1$}\\
\hline\hline
Exact & Numerics & Error \\
\hline
$0$ & $0$ & $0\%$\\ % I
$7/8$ & $0.8995$ & $+2.80\%$\\ % sigma'
$7/8+1$ & $1.9096$ & $+1.85\%$\\
$7/8+1$ & $1.9141$ & $+2.09\%$\\
$0+2$ & $2.0124$ & $+0.62\%$\\
$0+2$ & $2.0181$ & $+0.90\%$\\
\hline
\end{tabular}
~
\begin{tabular}{|c|c|c|}
\hline
\multicolumn{3}{|c|}{$y_1=y_2=\tau$}\\
\hline\hline
Exact & Numerics & Error \\
\hline
$3/40$ & $0.0751$ & $+0.19\%$\\ % sigma
$1/5$  & $0.2006$ & $+0.28\%$\\ % epsilon
$3/40+1$ & $1.0730$ & $-0.19\%$\\
$3/40+1$ & $1.0884$ & $+1.25\%$\\
$6/5$   & $1.2026$ & $+0.21\%$\\ % epsilon'
$1/5+1$ & $1.2156$ & $+1.30\%$\\
%1+1/5 & 1.2332 & $+2.77\%$\\
\hline
\end{tabular}

~

~

\begin{tabular}{|c|c|c|}
\hline
\multicolumn{3}{|c|}{$y_1=1$, $y_2=\tau$}\\
\hline\hline
Exact & Numerics & Error \\
\hline
$19/40$ & $0.4757$ & $+0.14\%$\\ % 3/40 + 2/5 : Spin 13/40
$3/5$   & $0.6009$ & $+0.15\%$\\ % 3/5 + 0    : Spin 24/40
19/40+1 & $1.4549$ & $-1.37\%$\\
19/40+1 & $1.5022$ & $+1.85\%$\\
$3/5+1$ & $1.5414$ & $-3.66\%$\\
$3/5+1$ & $1.6129$ & $+0.80\%$\\
\hline
\end{tabular}
~
\begin{tabular}{|c|c|c|}
\hline
\multicolumn{3}{|c|}{$y_1=\tau$, $y_2=1$}\\
\hline\hline
Exact & Numerics & Error \\
\hline
$19/40$ & $0.4757$ & $+0.14\%$\\
$3/5$   & $0.6009$ & $+0.15\%$\\
19/40+1 & $1.4549$ & $-1.37\%$\\
19/40+1 & $1.5022$ & $+1.85\%$\\
$3/5+1$ & $1.5414$ & $-3.66\%$\\
$3/5+1$ & $1.6129$ & $+0.80\%$\\
\hline
\end{tabular}
\caption{Scaling dimensions for Fibonacci anyons with antiferromagnetic nearest neigbour interactions on an infinite chain. Numerical values were computed using an anyonic MERA with maximum degeneracies for charges $1$ and $\tau$ of 3 and 5 respectively (denoted $\chi=[3,5]$), and are grouped according to their classification by the values of $y_1$ and $y_2$ in \protect{\fref{fig:scalingsuperop}}.
\label{tab:results}}
\end{table}
\begin{table}
\begin{tabular}{|c|c|c|}
\hline
\multicolumn{3}{|c|}{$y_1=y_2=1$}\\
\hline\hline
Exact & Numerics & Error \\
\hline
$0$ & $0$ & $0\%$\\ % I
$4/3$ & $1.3514$ & $+1.36\%$\\ % Z1
$4/3$ & $1.3695$ & $+2.71\%$\\ % Z2
$0+2$ & $1.9519$ & $-2.41\%$\\
$0+2$ & $1.9742$ & $-1.29\%$\\
$1+4/3$ & $2.2570$ & $-3.27\%$\\
\hline
\end{tabular}
~
\begin{tabular}{|c|c|c|}
\hline
\multicolumn{3}{|c|}{$y_1=y_2=\tau$}\\
\hline\hline
Exact & Numerics & Error \\
\hline
$2/15$ & $0.1329$ & $-0.35\%$\\ % sigma1
$2/15$  & $0.1339$ & $+0.44\%$\\ % sigma2
$4/5$ & $0.8134$ & $+1.67\%$\\ % epsilon
$2/15+1$ & $1.0937$ & $-3.49\%$\\
$2/15+1$   & $1.1108$ & $-1.99\%$\\
$2/15+1$ & $1.1622$ & $+2.55\%$\\
\hline
\end{tabular}

~

~

\begin{tabular}{|c|c|c|}
\hline
\multicolumn{3}{|c|}{$y_1=1$, $y_2=\tau$}\\
\hline\hline
Exact & Numerics & Error \\
\hline
$2/5$ & $0.3993$ & $-0.18\%$\\     % 2/5 + 0    : Spin 6/15
$11/15$   & $0.7327$ & $-0.09\%$\\ % 3/5 + 2/15 : Spin 7/15
$11/15$ & $0.7392$ & $+0.80\%$\\
$2/5+1$ & $1.3699$ & $-2.15\%$\\
$2/5+1$ & $1.3823$ & $-1.26\%$\\
11/15+1 & $1.6450$ & $-5.10\%$\\
\hline
\end{tabular}
~
\begin{tabular}{|c|c|c|}
\hline
\multicolumn{3}{|c|}{$y_1=\tau$, $y_2=1$}\\
\hline\hline
Exact & Numerics & Error \\
\hline
$2/5$ & $0.3993$ & $-0.18\%$\\
$11/15$   & $0.7327$ & $-0.09\%$\\
$11/15$ & $0.7392$ & $+0.80\%$\\
$2/5+1$ & $1.3699$ & $-2.15\%$\\
$2/5+1$ & $1.3823$ & $-1.26\%$\\
11/15+1 & $1.6450$ & $-5.10\%$\\
\hline
\end{tabular}
\caption{Scaling dimensions for Fibonacci anyons with ferromagnetic nearest neigbour interactions on an infinite chain. Numerical values were computed using an anyonic MERA with maximum degeneracies for charges $1$ and $\tau$ of 3 and 5 respectively (denoted $\chi=[3,5]$), and are grouped according to their classification by the values of $y_1$ and $y_2$ in \protect{\fref{fig:scalingsuperop}}. %We have omitted 
\label{tab:results2}}
\end{table}
\begin{figure}
%\includegraphics[width=246.0pt]{AFM}
%\\~\\
%\includegraphics[width=246.0pt]{FM}
\includegraphics[width=246.0pt]{graphs}
%\caption{Scaling dimensions %and conformal spins 
%for local scaling operators %for
%under (i) the AFM and (ii) the FM %interactions
%Hamiltonian on the golden chain. Results are grouped according to conformal towers.  The digits in the graph indicate the number of points within each grouping.
%%To determine conformal spin we exactly diagonalised the fixed point Hamiltonian on a four-site lattice (equivalent to 20 Fibonacci anyons), and computed the momenta of the energy eigenstates corresponding to the scaling operators of the conformal field theory. See also \rcite{feiguin2007,pfeifer2010a,difrancesco1996}.
\caption{COLOR ONLINE. Scaling dimensions of leading primary operators and their descendants, computed for %giving 
(i) antiferromagnetic and (ii) ferromagnetic local Hamiltonians \protect{\eref{eq:Hamiltonians}} on the golden chain. %nearest neighbour interactions.
Results are grouped into conformal towers, with a slight horizontal spread 
%allowing the reader to see that the MERA has correctly reproduced the %expected 
introduced to show the degeneracies of the descendant fields.
%digits on the graph indicating the number of eigenvalues within each group.
A circled cross indicates a primary field, and a plain cross indicates a descendant. Dashed lines indicate values predicted from CFT. 
%The conformal towers are labelled by the theoretical scaling dimensions of their primary operators and, where this is known, %there exists one, 
%%in the $(1,1)$ and $(\tau,\tau)$ sectors,
%the conventional name for the corresponding primary field. Operators arising from the %chiral 
%$(1,\tau)$ and $(\tau,1)$ sectors %of \protect{\fref{fig:scalingsuperop}} 
%are %nameless, and 
%identified %only 
%by their scaling dimension alone.
\label{fig:results}}
\end{figure}%


Comparison of the AFM case with existing results in the literature %for the antiferromagnetic chain 
show that the scaling dimensions obtained when $y_1=y_2$ 
correspond to those obtained when studying a system of anyons with a toroidal fusion diagram\citestop{feiguin2007} For a system of anyons on the torus it is possible to define an additional topological symmetry\pratext{ }\cite{feiguin2007} and classify local scaling operators according to whether or not they respect this symmetry. Operators satisfying $y_1=y_2=1$ correspond to those which respect the topological symmetry, and those satisfying $y_1=y_2=\tau$ do not. We will discuss the interpretation of the different sectors and their relationship to anyons on the torus in a forthcoming paper\citestop{pfeifer2010a}

When $y_1\not=y_2$ the scaling operators obtained are chiral, with those obtained from $y_1=1,~y_2=\tau$ and $y_1=\tau,~y_2=1$ believed to form conjugate pairs.

\section{Summary}

Numerical study of systems of interacting anyons is difficult due to their non-trivial exchange statistics. To date, study of these systems has been restricted to exact diagonalisation, Matrix Product States (MPS) for {1-D} systems, or special-case mappings to equivalent spin chains. This paper shows how any tensor network ansatz may be translated into a form applicable to systems of anyons, opening the door for the study of large systems of interacting anyons in both one and two dimensions. As an example, this paper demonstrates how the MERA may be implemented for a {1-D} anyonic system. This ansatz is particularly important as many
{1-D} systems of anyons are known which exhibit extended critical phases\pratext{ }\prbtext{.}\cite[see e.g.][]{feiguin2007,trebst2008,trebst2008a}\pratext{.} The structure of the MERA is known to be particularly well suited to reproducing long range correlations, and the scale-invariant MERA has the additional 
%further 
advantage of providing simple and direct means of computing the scaling dimensions and matrix representations of local scaling operators. %, and their scaling dimensions.

We applied the scale invariant MERA to infinite chains of Fibonacci anyons under antiferromagnetic and ferromagnetic nearest neighbour couplings, and identified a large number of local scaling operators. Our results for the scaling dimensions are in agreement with those previously obtained by exact diagonalisation of closely related systems, and for the relevant primary fields they are within $2.8\%$ of the theoretical values obtained from conformal field theory. 
We thus demonstrate that an anyonic MERA with $\chi=[3,5]$ permits conclusive identification of the relevant conformal field theory, and gives a level of accuracy comparable to that of the scale invariant MERA on a spin chain\citestop{pfeifer2009}

The anyonic generalisation of the {1-D} MERA presented here is useful in its own right, but the greatest significance of the approach described is that it is equally applicable to 2-D tensor network ans\"atze, and hence opens the door to studying the collective behaviour of large systems of anyons in two dimensions by numerical means, in situations where analytical solutions may not be possible.

The authors acknowledge the support of the Australian Research Council (FF0668731, DP0878830, DP1092513, APA). This research was supported in part by the Perimeter Institute for Theoretical Physics.

\emph{Note---}While preparing this paper for publication, we became aware of related work by \citeauthor{konig2010}\pratext{ }\prbtext{.}\cite{konig2010}\pratext{.} They also present the anyonic {1-D} MERA, providing proof of principle by
computing ground state energies for finite systems of Fibonacci anyons with $\chi=[1,1]$ ($s=2$ in their notation). 

%\appendix
%
%\section*{Appendix\label{sec:appendix}}
%
%This appendix lists expressions for certain tensors which are derivative of $(F^{abc}_d)_{(eu_1u_2)(fv_1v_2)}$. For 

% Create the reference section using BibTeX:
\bibliography{AnyonMERA}

\end{document}
%
% ****** End of file template.aps ******

